% Caja de Edgeworth del intercambio puro (3.A.21): gráfico maestro que crece por
% capas desde I.3 hasta III.2. Las variantes del PDF (caja-intercambio-equilibrio,
% caja-intercambio-teoremas…) fijan \cajacapas con las capas que muestran y
% cargan este fichero: mismo código y mismas cifras.
%
% Dotaciones totales X̄ = 10, Ȳ = 6; dotación ω = (1,5; 5) para A y (8,5; 1) para B.
% U_A = X_A^0,4 Y_A^0,6 (desde O_A); U_B = X_B^0,6 Y_B^0,4 (desde O_B).
% RMS^A_XY = (2/3) Y_A/X_A; RMS^B_XY = 1,5 Y_B/X_B. En ω: 2,22 y 0,18 (lente).
% Demandas (p_Y = 1, p = p_X): X_A = 0,4 (1,5p + 5)/p, Y_A = 0,6 (1,5p + 5);
%   X_B = 0,6 (8,5p + 1)/p, Y_B = 0,4 (8,5p + 1) (en la caja, X̄ - X_B, Ȳ - Y_B).
% Curvas de oferta-demanda: esos puntos al variar p.
% Equilibrio: p* = (0,4·5 + 0,6·1)/(0,6·1,5 + 0,4·8,5) = 0,6047;
%   E = (3,908; 3,544); U_A(E) = 3,685; U_B(E) = 4,236.
% Desequilibrio p = 0,23: A elige (9,296; 3,207) y B (2,291; 4,818) en la caja:
%   exceso de demanda de X = 7,004 y exceso de oferta de Y = 1,611 (0,23·7,004).
% Curva de contrato: Y_A = 9 X_A / ((2/3)(10 - X_A) + 1,5 X_A).
% Lente: U_A = U_A(ω) = 3,090 y U_B = U_B(ω) = 3,610; se cortan en X_A = 1,5 y 6,713.
% CC' (curva de contrato dentro de la lente): C = (3,148; 3,050), C' = (4,675; 3,983).
% 2TFEB: E' = (2,2; 2,329) en la curva de contrato, RMS = 0,706; la recta de
%   precios tangente en E' pasa por ω' = (1,5; 2,824): A cede 2,18 de Y a B.
% Autor: Víctor Gutiérrez Marcos
\documentclass[tikz,border=4pt]{standalone}
% ==========================================================================
% tikz-tcee.tex — Estilos comunes de los gráficos TikZ del temario
% Autor: Víctor Gutiérrez Marcos
%
% Cada gráfico es un fichero standalone en fuentes/<tema>/graficos/:
%
%   \documentclass[tikz,border=4pt]{standalone}
%   % ==========================================================================
% tikz-tcee.tex — Estilos comunes de los gráficos TikZ del temario
% Autor: Víctor Gutiérrez Marcos
%
% Cada gráfico es un fichero standalone en fuentes/<tema>/graficos/:
%
%   \documentclass[tikz,border=4pt]{standalone}
%   % ==========================================================================
% tikz-tcee.tex — Estilos comunes de los gráficos TikZ del temario
% Autor: Víctor Gutiérrez Marcos
%
% Cada gráfico es un fichero standalone en fuentes/<tema>/graficos/:
%
%   \documentclass[tikz,border=4pt]{standalone}
%   \input{../../../latex/tikz-tcee}
%   \begin{document}
%   \begin{tikzpicture}[tcee]
%     \ejes{Q}{P}{6}{5}
%     \draw[curva1] (0.5,4.5) -- (5,0.8) node[etiqueta,right] {$D$};
%     \capa{2}{\draw[curva2] ...;}   % en el vídeo aparece en el paso 2
%   \end{tikzpicture}
%   \end{document}
%
% Paleta: la de la web (morado, dorado, salvia) más tres colores de apoyo
% distinguibles en blanco y negro por el trazo.
% ==========================================================================
\usepackage[T1]{fontenc}
\usepackage{newpxtext,newpxmath}
\usepackage{xcolor}
\usepackage{tikz,pgfplots}
\pgfplotsset{compat=1.18}
\usetikzlibrary{arrows.meta,calc,positioning,patterns,decorations.pathreplacing,intersections,shapes.misc,backgrounds}

\definecolor{tceemorado}{HTML}{5F2987}
\definecolor{tceedorado}{HTML}{B8860B}
\definecolor{tceeverde}{HTML}{3C7D3A}
\definecolor{tceeazul}{HTML}{1F5F99}
\definecolor{tceerojo}{HTML}{B22222}
\definecolor{tceenaranja}{HTML}{E07B28}
\definecolor{tceegris}{HTML}{767171}
\definecolor{tceesalvia}{HTML}{E2EFD9}

% Capas para el vídeo: \capa{n}{…} dibuja su contenido solo a partir del paso n.
% En el PDF, la web y Word se ve todo; generar-video.py compila el gráfico una
% vez por paso con \def\tceepaso{k} para que cada elemento aparezca al narrarlo.
% Vectores con flecha encima, como en el texto (tcee.sty): \vect{x}, \vect{p}
\providecommand{\vect}[1]{\vec{#1}}
\providecommand{\tceepaso}{1000}
\newcommand{\capa}[2]{\ifnum#1>\tceepaso\relax\else#2\fi}

\tikzset{
  tcee/.style={line cap=round,line join=round,font=\small,>={Stealth[length=2.2mm]}},
  eje/.style={->,thick,black},
  curva1/.style={very thick,tceemorado},
  curva2/.style={very thick,tceeazul},
  curva3/.style={very thick,tceeverde},
  curva4/.style={very thick,tceedorado},
  curvanueva/.style={very thick,tceerojo},
  desplazada/.style={very thick,dashed},
  guia/.style={thin,dashed,tceegris},
  punto/.style={circle,fill=black,inner sep=1.3pt},
  etiqueta/.style={font=\small,inner sep=2pt},
  flecha/.style={->,thick,tceerojo},
  area/.style={fill=tceemorado,fill opacity=0.18,draw=none},
  area2/.style={fill=tceedorado,fill opacity=0.22,draw=none},
}

% \ejes{etiqueta x}{etiqueta y}{largo x}{alto y}
\newcommand{\ejes}[4]{%
  \draw[eje] (0,0) -- (#3,0) node[below left=2pt and 0pt] {#1};
  \draw[eje] (0,0) -- (0,#4) node[left] {#2};
  \node[below left] at (0,0) {$0$};}
% \proyeccion{(x,y)}{etq x}{etq y}: líneas guía a los ejes con etiquetas
\newcommand{\proyeccion}[3]{%
  \draw[guia] let \p1=#1 in (\x1,0) node[below] {#2} |- (0,\y1) node[left] {#3};}

% --------------------------------------------------------------------------
% Paleta «dictamen» (dic…): gráficos de datos de la Guía del dictamen (1D01) y
% de cualquier gráfico de coyuntura que quiera los mismos colores. Un concepto,
% un color en todos los gráficos (CRITERIOS §7, «tres reglas», 24/10/2026);
% tabla completa en fuentes/1D01/_trabajo/paleta-1D01.md. No cambia nada de
% lo anterior: son colores y estilos nuevos.
% Actividad: el PIB (y sus variantes: PIB por persona, PIB mundial), morado.
\definecolor{dicpib}{HTML}{5F2987}
% Demanda interna, gama cálida: consumo (naranjas: privado oscuro, público
% claro), inversión (verde) y existencias (gris con trama, aparte).
\definecolor{dicconspriv}{HTML}{C76A1F}
\definecolor{dicconspub}{HTML}{F5C690}
\definecolor{dicfbcf}{HTML}{2A6B3C}
\definecolor{dicexistencias}{HTML}{B4B4B4}
% Demanda externa y comercio, gama fría (azul claro, que se separa del morado)
\definecolor{dicexterna}{HTML}{6AA6DA}
% Zonas: zona del euro (azul marino), España (rojo), resto de países (grises)
\definecolor{dicze}{HTML}{17406D}
\definecolor{dicespana}{HTML}{C8382E}
\definecolor{dicotro}{HTML}{8A8A8A}
\definecolor{dicotroclaro}{HTML}{B5B5B5}
\definecolor{dicotromasclaro}{HTML}{DADADA}
% Precios: inflación general (negro, continuo) y subyacente (negro,
% discontinuo); componentes: energía (bermellón), alimentos (amarillo) y los
% de la subyacente en azules (bienes industriales claro, servicios oscuro).
\definecolor{dicinflacion}{HTML}{1A1A1A}
\definecolor{dicenergia}{HTML}{D55E00}
\definecolor{dicalimentos}{HTML}{F2D45C}
\definecolor{dicbine}{HTML}{7FB2E0}
\definecolor{dicservicios}{HTML}{1F4E89}
% Trabajo: empleo (verde azulado), horas (el mismo, claro), tasa de actividad
% (verde azulado oscuro), paro (teja); productividad (dorado de la web).
\definecolor{dicempleo}{HTML}{1B7F6B}
\definecolor{dichoras}{HTML}{86CBB8}
\definecolor{dicactividad}{HTML}{0F5247}
\definecolor{dicparo}{HTML}{A23B2A}
\definecolor{dicproductividad}{HTML}{B8860B}
% Hogares: renta disponible (pizarra) y tasa de ahorro (negro); el consumo es
% dicconspriv.
\definecolor{dicrenta}{HTML}{4A5A6A}
\definecolor{dicahorro}{HTML}{1A1A1A}
% Política monetaria: tipos del BCE en azules (depósito, el de la zona del
% euro; operaciones principales, medio; facilidad marginal, claro); balance:
% préstamos a los bancos (dorado claro), resto del activo (gris claro).
\definecolor{dicdeposito}{HTML}{17406D}
\definecolor{dicopf}{HTML}{4F86C6}
\definecolor{dicfmc}{HTML}{9EC5EA}
\definecolor{dicprestamos}{HTML}{D9A520}
\definecolor{dicresto}{HTML}{CFCFCF}
% Ciclo: series adelantadas, coincidentes (el PIB) y retrasadas
\definecolor{dicadelantada}{HTML}{D55E00}
\definecolor{dicretrasada}{HTML}{6AA6DA}
% Referencias: rejilla, previsión, sombreados
\definecolor{dicreferencia}{HTML}{767171}

\tikzset{
  % barra de un componente: \addplot[dicbarra=dicfbcf]
  dicbarra/.style={fill=#1, draw=#1!65!black, line width=0.3pt},
  % existencias: gris con trama, para que se distinga también en gris
  dicbarratrama/.style={fill=#1, draw=#1!65!black, line width=0.3pt,
    postaction={pattern=north east lines, pattern color=white}},
  % línea de una serie; con halo blanco cuando pasa sobre barras
  diclinea/.style={very thick, #1},
  diclineahalo/.style={very thick, #1, preaction={draw=white, line width=3.4pt, solid}},
  % puntos de una serie (capa aparte, después de la línea)
  dicpuntos/.style={only marks, mark=*, mark size=1.8pt,
    mark options={fill=white, draw=#1, line width=0.9pt, solid}},
  % línea del cero o de un umbral, después de las barras
  diccero/.style={thin, black},
  dicumbral/.style={thin, dashed, dicreferencia},
}
% Capas de los gráficos de datos (CRITERIOS §7, «tres reglas»), de abajo
% arriba: cuadrícula → sombreados (dicfondo) → barras y áreas (dicbarras) → ejes, cero y umbrales
% (axis lines) → líneas (main) → puntos (dicpuntos) → rótulos (axis
% descriptions). No basta el orden del código: pgfplots dibuja cuadrícula,
% ejes y marcas en sus propias capas (y area style pone los ejes encima), y
% dentro de un axis los \draw, \fill y \node no admiten capa (van con los
% datos de main); por eso todo lo que se dibuja dentro del axis es un \addplot.
% Uso: \begin{axis}[…, dicapas] (al final de las opciones);
%   barras y áreas: \addplot[on layer=dicbarras, …];
%   sombreados:     \dicsombra[estilo]{x1}{y1}{x2}{y2} (rectángulo, en dicfondo,
%                   bajo las barras y sobre la cuadrícula);
%   cero y umbrales: \dicraya[estilo]{(x1,y1) (x2,y2)} (en axis lines);
%   líneas: \addplot normal (main); puntos: \addplot[on layer=dicpuntos, …]
%   (las marcas de cualquier serie van solas a dicpuntos);
%   rótulos: \dicrotulo[opciones del nodo]{(x,y)}{texto} (en axis descriptions),
%   siempre sin fondo (nada de fill=white: si chocan con una línea, se desplazan;
%   CRITERIOS §7, regla 4).
% Los \capa{n} del vídeo no cambian: siguen envolviendo instrucciones completas.
\pgfdeclarelayer{dicfondo}
\pgfdeclarelayer{dicbarras}
\pgfdeclarelayer{dicpuntos}
\pgfplotsset{
  % (el segundo argumento reparte cuadrícula, ejes y rótulos como el estándar)
  layers/dic/.define layer set={axis background, axis grid, dicfondo, dicbarras,
    axis ticks, axis lines, axis tick labels, pre main, main, dicpuntos,
    axis descriptions, axis foreground}{/pgfplots/layers/standard},
  % leyenda sin caja blanca: toma el fondo del gráfico, también si el gráfico va
  % dentro de un cuadro gris (CRITERIOS §7, regla 4, 26/10/2026)
  dicapas/.style={axis on top=false, set layers=dic, mark layer=dicpuntos,
    legend style={fill=none}},
  dicsuelto/.style={forget plot, stack plots=false, sharp plot, mark=none},
}
\newcommand{\dicraya}[2][]{%
  \addplot[on layer=axis lines, dicsuelto, diccero, #1] coordinates {#2};}
\newcommand{\dicsombra}[5][]{%
  \addplot[on layer=dicfondo, dicsuelto, draw=none, #1]
    coordinates {(#2,#3) (#4,#3) (#4,#5) (#2,#5)};}
\newcommand{\dicrotulo}[3][]{%
  \addplot[on layer=axis descriptions, dicsuelto, draw=none] coordinates {#2}
    node[pos=0, #1] {#3};}

%   \begin{document}
%   \begin{tikzpicture}[tcee]
%     \ejes{Q}{P}{6}{5}
%     \draw[curva1] (0.5,4.5) -- (5,0.8) node[etiqueta,right] {$D$};
%     \capa{2}{\draw[curva2] ...;}   % en el vídeo aparece en el paso 2
%   \end{tikzpicture}
%   \end{document}
%
% Paleta: la de la web (morado, dorado, salvia) más tres colores de apoyo
% distinguibles en blanco y negro por el trazo.
% ==========================================================================
\usepackage[T1]{fontenc}
\usepackage{newpxtext,newpxmath}
\usepackage{xcolor}
\usepackage{tikz,pgfplots}
\pgfplotsset{compat=1.18}
\usetikzlibrary{arrows.meta,calc,positioning,patterns,decorations.pathreplacing,intersections,shapes.misc,backgrounds}

\definecolor{tceemorado}{HTML}{5F2987}
\definecolor{tceedorado}{HTML}{B8860B}
\definecolor{tceeverde}{HTML}{3C7D3A}
\definecolor{tceeazul}{HTML}{1F5F99}
\definecolor{tceerojo}{HTML}{B22222}
\definecolor{tceenaranja}{HTML}{E07B28}
\definecolor{tceegris}{HTML}{767171}
\definecolor{tceesalvia}{HTML}{E2EFD9}

% Capas para el vídeo: \capa{n}{…} dibuja su contenido solo a partir del paso n.
% En el PDF, la web y Word se ve todo; generar-video.py compila el gráfico una
% vez por paso con \def\tceepaso{k} para que cada elemento aparezca al narrarlo.
% Vectores con flecha encima, como en el texto (tcee.sty): \vect{x}, \vect{p}
\providecommand{\vect}[1]{\vec{#1}}
\providecommand{\tceepaso}{1000}
\newcommand{\capa}[2]{\ifnum#1>\tceepaso\relax\else#2\fi}

\tikzset{
  tcee/.style={line cap=round,line join=round,font=\small,>={Stealth[length=2.2mm]}},
  eje/.style={->,thick,black},
  curva1/.style={very thick,tceemorado},
  curva2/.style={very thick,tceeazul},
  curva3/.style={very thick,tceeverde},
  curva4/.style={very thick,tceedorado},
  curvanueva/.style={very thick,tceerojo},
  desplazada/.style={very thick,dashed},
  guia/.style={thin,dashed,tceegris},
  punto/.style={circle,fill=black,inner sep=1.3pt},
  etiqueta/.style={font=\small,inner sep=2pt},
  flecha/.style={->,thick,tceerojo},
  area/.style={fill=tceemorado,fill opacity=0.18,draw=none},
  area2/.style={fill=tceedorado,fill opacity=0.22,draw=none},
}

% \ejes{etiqueta x}{etiqueta y}{largo x}{alto y}
\newcommand{\ejes}[4]{%
  \draw[eje] (0,0) -- (#3,0) node[below left=2pt and 0pt] {#1};
  \draw[eje] (0,0) -- (0,#4) node[left] {#2};
  \node[below left] at (0,0) {$0$};}
% \proyeccion{(x,y)}{etq x}{etq y}: líneas guía a los ejes con etiquetas
\newcommand{\proyeccion}[3]{%
  \draw[guia] let \p1=#1 in (\x1,0) node[below] {#2} |- (0,\y1) node[left] {#3};}

% --------------------------------------------------------------------------
% Paleta «dictamen» (dic…): gráficos de datos de la Guía del dictamen (1D01) y
% de cualquier gráfico de coyuntura que quiera los mismos colores. Un concepto,
% un color en todos los gráficos (CRITERIOS §7, «tres reglas», 24/10/2026);
% tabla completa en fuentes/1D01/_trabajo/paleta-1D01.md. No cambia nada de
% lo anterior: son colores y estilos nuevos.
% Actividad: el PIB (y sus variantes: PIB por persona, PIB mundial), morado.
\definecolor{dicpib}{HTML}{5F2987}
% Demanda interna, gama cálida: consumo (naranjas: privado oscuro, público
% claro), inversión (verde) y existencias (gris con trama, aparte).
\definecolor{dicconspriv}{HTML}{C76A1F}
\definecolor{dicconspub}{HTML}{F5C690}
\definecolor{dicfbcf}{HTML}{2A6B3C}
\definecolor{dicexistencias}{HTML}{B4B4B4}
% Demanda externa y comercio, gama fría (azul claro, que se separa del morado)
\definecolor{dicexterna}{HTML}{6AA6DA}
% Zonas: zona del euro (azul marino), España (rojo), resto de países (grises)
\definecolor{dicze}{HTML}{17406D}
\definecolor{dicespana}{HTML}{C8382E}
\definecolor{dicotro}{HTML}{8A8A8A}
\definecolor{dicotroclaro}{HTML}{B5B5B5}
\definecolor{dicotromasclaro}{HTML}{DADADA}
% Precios: inflación general (negro, continuo) y subyacente (negro,
% discontinuo); componentes: energía (bermellón), alimentos (amarillo) y los
% de la subyacente en azules (bienes industriales claro, servicios oscuro).
\definecolor{dicinflacion}{HTML}{1A1A1A}
\definecolor{dicenergia}{HTML}{D55E00}
\definecolor{dicalimentos}{HTML}{F2D45C}
\definecolor{dicbine}{HTML}{7FB2E0}
\definecolor{dicservicios}{HTML}{1F4E89}
% Trabajo: empleo (verde azulado), horas (el mismo, claro), tasa de actividad
% (verde azulado oscuro), paro (teja); productividad (dorado de la web).
\definecolor{dicempleo}{HTML}{1B7F6B}
\definecolor{dichoras}{HTML}{86CBB8}
\definecolor{dicactividad}{HTML}{0F5247}
\definecolor{dicparo}{HTML}{A23B2A}
\definecolor{dicproductividad}{HTML}{B8860B}
% Hogares: renta disponible (pizarra) y tasa de ahorro (negro); el consumo es
% dicconspriv.
\definecolor{dicrenta}{HTML}{4A5A6A}
\definecolor{dicahorro}{HTML}{1A1A1A}
% Política monetaria: tipos del BCE en azules (depósito, el de la zona del
% euro; operaciones principales, medio; facilidad marginal, claro); balance:
% préstamos a los bancos (dorado claro), resto del activo (gris claro).
\definecolor{dicdeposito}{HTML}{17406D}
\definecolor{dicopf}{HTML}{4F86C6}
\definecolor{dicfmc}{HTML}{9EC5EA}
\definecolor{dicprestamos}{HTML}{D9A520}
\definecolor{dicresto}{HTML}{CFCFCF}
% Ciclo: series adelantadas, coincidentes (el PIB) y retrasadas
\definecolor{dicadelantada}{HTML}{D55E00}
\definecolor{dicretrasada}{HTML}{6AA6DA}
% Referencias: rejilla, previsión, sombreados
\definecolor{dicreferencia}{HTML}{767171}

\tikzset{
  % barra de un componente: \addplot[dicbarra=dicfbcf]
  dicbarra/.style={fill=#1, draw=#1!65!black, line width=0.3pt},
  % existencias: gris con trama, para que se distinga también en gris
  dicbarratrama/.style={fill=#1, draw=#1!65!black, line width=0.3pt,
    postaction={pattern=north east lines, pattern color=white}},
  % línea de una serie; con halo blanco cuando pasa sobre barras
  diclinea/.style={very thick, #1},
  diclineahalo/.style={very thick, #1, preaction={draw=white, line width=3.4pt, solid}},
  % puntos de una serie (capa aparte, después de la línea)
  dicpuntos/.style={only marks, mark=*, mark size=1.8pt,
    mark options={fill=white, draw=#1, line width=0.9pt, solid}},
  % línea del cero o de un umbral, después de las barras
  diccero/.style={thin, black},
  dicumbral/.style={thin, dashed, dicreferencia},
}
% Capas de los gráficos de datos (CRITERIOS §7, «tres reglas»), de abajo
% arriba: cuadrícula → sombreados (dicfondo) → barras y áreas (dicbarras) → ejes, cero y umbrales
% (axis lines) → líneas (main) → puntos (dicpuntos) → rótulos (axis
% descriptions). No basta el orden del código: pgfplots dibuja cuadrícula,
% ejes y marcas en sus propias capas (y area style pone los ejes encima), y
% dentro de un axis los \draw, \fill y \node no admiten capa (van con los
% datos de main); por eso todo lo que se dibuja dentro del axis es un \addplot.
% Uso: \begin{axis}[…, dicapas] (al final de las opciones);
%   barras y áreas: \addplot[on layer=dicbarras, …];
%   sombreados:     \dicsombra[estilo]{x1}{y1}{x2}{y2} (rectángulo, en dicfondo,
%                   bajo las barras y sobre la cuadrícula);
%   cero y umbrales: \dicraya[estilo]{(x1,y1) (x2,y2)} (en axis lines);
%   líneas: \addplot normal (main); puntos: \addplot[on layer=dicpuntos, …]
%   (las marcas de cualquier serie van solas a dicpuntos);
%   rótulos: \dicrotulo[opciones del nodo]{(x,y)}{texto} (en axis descriptions),
%   siempre sin fondo (nada de fill=white: si chocan con una línea, se desplazan;
%   CRITERIOS §7, regla 4).
% Los \capa{n} del vídeo no cambian: siguen envolviendo instrucciones completas.
\pgfdeclarelayer{dicfondo}
\pgfdeclarelayer{dicbarras}
\pgfdeclarelayer{dicpuntos}
\pgfplotsset{
  % (el segundo argumento reparte cuadrícula, ejes y rótulos como el estándar)
  layers/dic/.define layer set={axis background, axis grid, dicfondo, dicbarras,
    axis ticks, axis lines, axis tick labels, pre main, main, dicpuntos,
    axis descriptions, axis foreground}{/pgfplots/layers/standard},
  % leyenda sin caja blanca: toma el fondo del gráfico, también si el gráfico va
  % dentro de un cuadro gris (CRITERIOS §7, regla 4, 26/10/2026)
  dicapas/.style={axis on top=false, set layers=dic, mark layer=dicpuntos,
    legend style={fill=none}},
  dicsuelto/.style={forget plot, stack plots=false, sharp plot, mark=none},
}
\newcommand{\dicraya}[2][]{%
  \addplot[on layer=axis lines, dicsuelto, diccero, #1] coordinates {#2};}
\newcommand{\dicsombra}[5][]{%
  \addplot[on layer=dicfondo, dicsuelto, draw=none, #1]
    coordinates {(#2,#3) (#4,#3) (#4,#5) (#2,#5)};}
\newcommand{\dicrotulo}[3][]{%
  \addplot[on layer=axis descriptions, dicsuelto, draw=none] coordinates {#2}
    node[pos=0, #1] {#3};}

%   \begin{document}
%   \begin{tikzpicture}[tcee]
%     \ejes{Q}{P}{6}{5}
%     \draw[curva1] (0.5,4.5) -- (5,0.8) node[etiqueta,right] {$D$};
%     \capa{2}{\draw[curva2] ...;}   % en el vídeo aparece en el paso 2
%   \end{tikzpicture}
%   \end{document}
%
% Paleta: la de la web (morado, dorado, salvia) más tres colores de apoyo
% distinguibles en blanco y negro por el trazo.
% ==========================================================================
\usepackage[T1]{fontenc}
\usepackage{newpxtext,newpxmath}
\usepackage{xcolor}
\usepackage{tikz,pgfplots}
\pgfplotsset{compat=1.18}
\usetikzlibrary{arrows.meta,calc,positioning,patterns,decorations.pathreplacing,intersections,shapes.misc,backgrounds}

\definecolor{tceemorado}{HTML}{5F2987}
\definecolor{tceedorado}{HTML}{B8860B}
\definecolor{tceeverde}{HTML}{3C7D3A}
\definecolor{tceeazul}{HTML}{1F5F99}
\definecolor{tceerojo}{HTML}{B22222}
\definecolor{tceenaranja}{HTML}{E07B28}
\definecolor{tceegris}{HTML}{767171}
\definecolor{tceesalvia}{HTML}{E2EFD9}

% Capas para el vídeo: \capa{n}{…} dibuja su contenido solo a partir del paso n.
% En el PDF, la web y Word se ve todo; generar-video.py compila el gráfico una
% vez por paso con \def\tceepaso{k} para que cada elemento aparezca al narrarlo.
% Vectores con flecha encima, como en el texto (tcee.sty): \vect{x}, \vect{p}
\providecommand{\vect}[1]{\vec{#1}}
\providecommand{\tceepaso}{1000}
\newcommand{\capa}[2]{\ifnum#1>\tceepaso\relax\else#2\fi}

\tikzset{
  tcee/.style={line cap=round,line join=round,font=\small,>={Stealth[length=2.2mm]}},
  eje/.style={->,thick,black},
  curva1/.style={very thick,tceemorado},
  curva2/.style={very thick,tceeazul},
  curva3/.style={very thick,tceeverde},
  curva4/.style={very thick,tceedorado},
  curvanueva/.style={very thick,tceerojo},
  desplazada/.style={very thick,dashed},
  guia/.style={thin,dashed,tceegris},
  punto/.style={circle,fill=black,inner sep=1.3pt},
  etiqueta/.style={font=\small,inner sep=2pt},
  flecha/.style={->,thick,tceerojo},
  area/.style={fill=tceemorado,fill opacity=0.18,draw=none},
  area2/.style={fill=tceedorado,fill opacity=0.22,draw=none},
}

% \ejes{etiqueta x}{etiqueta y}{largo x}{alto y}
\newcommand{\ejes}[4]{%
  \draw[eje] (0,0) -- (#3,0) node[below left=2pt and 0pt] {#1};
  \draw[eje] (0,0) -- (0,#4) node[left] {#2};
  \node[below left] at (0,0) {$0$};}
% \proyeccion{(x,y)}{etq x}{etq y}: líneas guía a los ejes con etiquetas
\newcommand{\proyeccion}[3]{%
  \draw[guia] let \p1=#1 in (\x1,0) node[below] {#2} |- (0,\y1) node[left] {#3};}

% --------------------------------------------------------------------------
% Paleta «dictamen» (dic…): gráficos de datos de la Guía del dictamen (1D01) y
% de cualquier gráfico de coyuntura que quiera los mismos colores. Un concepto,
% un color en todos los gráficos (CRITERIOS §7, «tres reglas», 24/10/2026);
% tabla completa en fuentes/1D01/_trabajo/paleta-1D01.md. No cambia nada de
% lo anterior: son colores y estilos nuevos.
% Actividad: el PIB (y sus variantes: PIB por persona, PIB mundial), morado.
\definecolor{dicpib}{HTML}{5F2987}
% Demanda interna, gama cálida: consumo (naranjas: privado oscuro, público
% claro), inversión (verde) y existencias (gris con trama, aparte).
\definecolor{dicconspriv}{HTML}{C76A1F}
\definecolor{dicconspub}{HTML}{F5C690}
\definecolor{dicfbcf}{HTML}{2A6B3C}
\definecolor{dicexistencias}{HTML}{B4B4B4}
% Demanda externa y comercio, gama fría (azul claro, que se separa del morado)
\definecolor{dicexterna}{HTML}{6AA6DA}
% Zonas: zona del euro (azul marino), España (rojo), resto de países (grises)
\definecolor{dicze}{HTML}{17406D}
\definecolor{dicespana}{HTML}{C8382E}
\definecolor{dicotro}{HTML}{8A8A8A}
\definecolor{dicotroclaro}{HTML}{B5B5B5}
\definecolor{dicotromasclaro}{HTML}{DADADA}
% Precios: inflación general (negro, continuo) y subyacente (negro,
% discontinuo); componentes: energía (bermellón), alimentos (amarillo) y los
% de la subyacente en azules (bienes industriales claro, servicios oscuro).
\definecolor{dicinflacion}{HTML}{1A1A1A}
\definecolor{dicenergia}{HTML}{D55E00}
\definecolor{dicalimentos}{HTML}{F2D45C}
\definecolor{dicbine}{HTML}{7FB2E0}
\definecolor{dicservicios}{HTML}{1F4E89}
% Trabajo: empleo (verde azulado), horas (el mismo, claro), tasa de actividad
% (verde azulado oscuro), paro (teja); productividad (dorado de la web).
\definecolor{dicempleo}{HTML}{1B7F6B}
\definecolor{dichoras}{HTML}{86CBB8}
\definecolor{dicactividad}{HTML}{0F5247}
\definecolor{dicparo}{HTML}{A23B2A}
\definecolor{dicproductividad}{HTML}{B8860B}
% Hogares: renta disponible (pizarra) y tasa de ahorro (negro); el consumo es
% dicconspriv.
\definecolor{dicrenta}{HTML}{4A5A6A}
\definecolor{dicahorro}{HTML}{1A1A1A}
% Política monetaria: tipos del BCE en azules (depósito, el de la zona del
% euro; operaciones principales, medio; facilidad marginal, claro); balance:
% préstamos a los bancos (dorado claro), resto del activo (gris claro).
\definecolor{dicdeposito}{HTML}{17406D}
\definecolor{dicopf}{HTML}{4F86C6}
\definecolor{dicfmc}{HTML}{9EC5EA}
\definecolor{dicprestamos}{HTML}{D9A520}
\definecolor{dicresto}{HTML}{CFCFCF}
% Ciclo: series adelantadas, coincidentes (el PIB) y retrasadas
\definecolor{dicadelantada}{HTML}{D55E00}
\definecolor{dicretrasada}{HTML}{6AA6DA}
% Referencias: rejilla, previsión, sombreados
\definecolor{dicreferencia}{HTML}{767171}

\tikzset{
  % barra de un componente: \addplot[dicbarra=dicfbcf]
  dicbarra/.style={fill=#1, draw=#1!65!black, line width=0.3pt},
  % existencias: gris con trama, para que se distinga también en gris
  dicbarratrama/.style={fill=#1, draw=#1!65!black, line width=0.3pt,
    postaction={pattern=north east lines, pattern color=white}},
  % línea de una serie; con halo blanco cuando pasa sobre barras
  diclinea/.style={very thick, #1},
  diclineahalo/.style={very thick, #1, preaction={draw=white, line width=3.4pt, solid}},
  % puntos de una serie (capa aparte, después de la línea)
  dicpuntos/.style={only marks, mark=*, mark size=1.8pt,
    mark options={fill=white, draw=#1, line width=0.9pt, solid}},
  % línea del cero o de un umbral, después de las barras
  diccero/.style={thin, black},
  dicumbral/.style={thin, dashed, dicreferencia},
}
% Capas de los gráficos de datos (CRITERIOS §7, «tres reglas»), de abajo
% arriba: cuadrícula → sombreados (dicfondo) → barras y áreas (dicbarras) → ejes, cero y umbrales
% (axis lines) → líneas (main) → puntos (dicpuntos) → rótulos (axis
% descriptions). No basta el orden del código: pgfplots dibuja cuadrícula,
% ejes y marcas en sus propias capas (y area style pone los ejes encima), y
% dentro de un axis los \draw, \fill y \node no admiten capa (van con los
% datos de main); por eso todo lo que se dibuja dentro del axis es un \addplot.
% Uso: \begin{axis}[…, dicapas] (al final de las opciones);
%   barras y áreas: \addplot[on layer=dicbarras, …];
%   sombreados:     \dicsombra[estilo]{x1}{y1}{x2}{y2} (rectángulo, en dicfondo,
%                   bajo las barras y sobre la cuadrícula);
%   cero y umbrales: \dicraya[estilo]{(x1,y1) (x2,y2)} (en axis lines);
%   líneas: \addplot normal (main); puntos: \addplot[on layer=dicpuntos, …]
%   (las marcas de cualquier serie van solas a dicpuntos);
%   rótulos: \dicrotulo[opciones del nodo]{(x,y)}{texto} (en axis descriptions),
%   siempre sin fondo (nada de fill=white: si chocan con una línea, se desplazan;
%   CRITERIOS §7, regla 4).
% Los \capa{n} del vídeo no cambian: siguen envolviendo instrucciones completas.
\pgfdeclarelayer{dicfondo}
\pgfdeclarelayer{dicbarras}
\pgfdeclarelayer{dicpuntos}
\pgfplotsset{
  % (el segundo argumento reparte cuadrícula, ejes y rótulos como el estándar)
  layers/dic/.define layer set={axis background, axis grid, dicfondo, dicbarras,
    axis ticks, axis lines, axis tick labels, pre main, main, dicpuntos,
    axis descriptions, axis foreground}{/pgfplots/layers/standard},
  % leyenda sin caja blanca: toma el fondo del gráfico, también si el gráfico va
  % dentro de un cuadro gris (CRITERIOS §7, regla 4, 26/10/2026)
  dicapas/.style={axis on top=false, set layers=dic, mark layer=dicpuntos,
    legend style={fill=none}},
  dicsuelto/.style={forget plot, stack plots=false, sharp plot, mark=none},
}
\newcommand{\dicraya}[2][]{%
  \addplot[on layer=axis lines, dicsuelto, diccero, #1] coordinates {#2};}
\newcommand{\dicsombra}[5][]{%
  \addplot[on layer=dicfondo, dicsuelto, draw=none, #1]
    coordinates {(#2,#3) (#4,#3) (#4,#5) (#2,#5)};}
\newcommand{\dicrotulo}[3][]{%
  \addplot[on layer=axis descriptions, dicsuelto, draw=none] coordinates {#2}
    node[pos=0, #1] {#3};}

% Variantes: \def\cajacapas{1,2,6} antes de % Caja de Edgeworth del intercambio puro (3.A.21): gráfico maestro que crece por
% capas desde I.3 hasta III.2. Las variantes del PDF (caja-intercambio-equilibrio,
% caja-intercambio-teoremas…) fijan \cajacapas con las capas que muestran y
% cargan este fichero: mismo código y mismas cifras.
%
% Dotaciones totales X̄ = 10, Ȳ = 6; dotación ω = (1,5; 5) para A y (8,5; 1) para B.
% U_A = X_A^0,4 Y_A^0,6 (desde O_A); U_B = X_B^0,6 Y_B^0,4 (desde O_B).
% RMS^A_XY = (2/3) Y_A/X_A; RMS^B_XY = 1,5 Y_B/X_B. En ω: 2,22 y 0,18 (lente).
% Demandas (p_Y = 1, p = p_X): X_A = 0,4 (1,5p + 5)/p, Y_A = 0,6 (1,5p + 5);
%   X_B = 0,6 (8,5p + 1)/p, Y_B = 0,4 (8,5p + 1) (en la caja, X̄ - X_B, Ȳ - Y_B).
% Curvas de oferta-demanda: esos puntos al variar p.
% Equilibrio: p* = (0,4·5 + 0,6·1)/(0,6·1,5 + 0,4·8,5) = 0,6047;
%   E = (3,908; 3,544); U_A(E) = 3,685; U_B(E) = 4,236.
% Desequilibrio p = 0,23: A elige (9,296; 3,207) y B (2,291; 4,818) en la caja:
%   exceso de demanda de X = 7,004 y exceso de oferta de Y = 1,611 (0,23·7,004).
% Curva de contrato: Y_A = 9 X_A / ((2/3)(10 - X_A) + 1,5 X_A).
% Lente: U_A = U_A(ω) = 3,090 y U_B = U_B(ω) = 3,610; se cortan en X_A = 1,5 y 6,713.
% CC' (curva de contrato dentro de la lente): C = (3,148; 3,050), C' = (4,675; 3,983).
% 2TFEB: E' = (2,2; 2,329) en la curva de contrato, RMS = 0,706; la recta de
%   precios tangente en E' pasa por ω' = (1,5; 2,824): A cede 2,18 de Y a B.
% Autor: Víctor Gutiérrez Marcos
\documentclass[tikz,border=4pt]{standalone}
% ==========================================================================
% tikz-tcee.tex — Estilos comunes de los gráficos TikZ del temario
% Autor: Víctor Gutiérrez Marcos
%
% Cada gráfico es un fichero standalone en fuentes/<tema>/graficos/:
%
%   \documentclass[tikz,border=4pt]{standalone}
%   % ==========================================================================
% tikz-tcee.tex — Estilos comunes de los gráficos TikZ del temario
% Autor: Víctor Gutiérrez Marcos
%
% Cada gráfico es un fichero standalone en fuentes/<tema>/graficos/:
%
%   \documentclass[tikz,border=4pt]{standalone}
%   \input{../../../latex/tikz-tcee}
%   \begin{document}
%   \begin{tikzpicture}[tcee]
%     \ejes{Q}{P}{6}{5}
%     \draw[curva1] (0.5,4.5) -- (5,0.8) node[etiqueta,right] {$D$};
%     \capa{2}{\draw[curva2] ...;}   % en el vídeo aparece en el paso 2
%   \end{tikzpicture}
%   \end{document}
%
% Paleta: la de la web (morado, dorado, salvia) más tres colores de apoyo
% distinguibles en blanco y negro por el trazo.
% ==========================================================================
\usepackage[T1]{fontenc}
\usepackage{newpxtext,newpxmath}
\usepackage{xcolor}
\usepackage{tikz,pgfplots}
\pgfplotsset{compat=1.18}
\usetikzlibrary{arrows.meta,calc,positioning,patterns,decorations.pathreplacing,intersections,shapes.misc,backgrounds}

\definecolor{tceemorado}{HTML}{5F2987}
\definecolor{tceedorado}{HTML}{B8860B}
\definecolor{tceeverde}{HTML}{3C7D3A}
\definecolor{tceeazul}{HTML}{1F5F99}
\definecolor{tceerojo}{HTML}{B22222}
\definecolor{tceenaranja}{HTML}{E07B28}
\definecolor{tceegris}{HTML}{767171}
\definecolor{tceesalvia}{HTML}{E2EFD9}

% Capas para el vídeo: \capa{n}{…} dibuja su contenido solo a partir del paso n.
% En el PDF, la web y Word se ve todo; generar-video.py compila el gráfico una
% vez por paso con \def\tceepaso{k} para que cada elemento aparezca al narrarlo.
% Vectores con flecha encima, como en el texto (tcee.sty): \vect{x}, \vect{p}
\providecommand{\vect}[1]{\vec{#1}}
\providecommand{\tceepaso}{1000}
\newcommand{\capa}[2]{\ifnum#1>\tceepaso\relax\else#2\fi}

\tikzset{
  tcee/.style={line cap=round,line join=round,font=\small,>={Stealth[length=2.2mm]}},
  eje/.style={->,thick,black},
  curva1/.style={very thick,tceemorado},
  curva2/.style={very thick,tceeazul},
  curva3/.style={very thick,tceeverde},
  curva4/.style={very thick,tceedorado},
  curvanueva/.style={very thick,tceerojo},
  desplazada/.style={very thick,dashed},
  guia/.style={thin,dashed,tceegris},
  punto/.style={circle,fill=black,inner sep=1.3pt},
  etiqueta/.style={font=\small,inner sep=2pt},
  flecha/.style={->,thick,tceerojo},
  area/.style={fill=tceemorado,fill opacity=0.18,draw=none},
  area2/.style={fill=tceedorado,fill opacity=0.22,draw=none},
}

% \ejes{etiqueta x}{etiqueta y}{largo x}{alto y}
\newcommand{\ejes}[4]{%
  \draw[eje] (0,0) -- (#3,0) node[below left=2pt and 0pt] {#1};
  \draw[eje] (0,0) -- (0,#4) node[left] {#2};
  \node[below left] at (0,0) {$0$};}
% \proyeccion{(x,y)}{etq x}{etq y}: líneas guía a los ejes con etiquetas
\newcommand{\proyeccion}[3]{%
  \draw[guia] let \p1=#1 in (\x1,0) node[below] {#2} |- (0,\y1) node[left] {#3};}

% --------------------------------------------------------------------------
% Paleta «dictamen» (dic…): gráficos de datos de la Guía del dictamen (1D01) y
% de cualquier gráfico de coyuntura que quiera los mismos colores. Un concepto,
% un color en todos los gráficos (CRITERIOS §7, «tres reglas», 24/10/2026);
% tabla completa en fuentes/1D01/_trabajo/paleta-1D01.md. No cambia nada de
% lo anterior: son colores y estilos nuevos.
% Actividad: el PIB (y sus variantes: PIB por persona, PIB mundial), morado.
\definecolor{dicpib}{HTML}{5F2987}
% Demanda interna, gama cálida: consumo (naranjas: privado oscuro, público
% claro), inversión (verde) y existencias (gris con trama, aparte).
\definecolor{dicconspriv}{HTML}{C76A1F}
\definecolor{dicconspub}{HTML}{F5C690}
\definecolor{dicfbcf}{HTML}{2A6B3C}
\definecolor{dicexistencias}{HTML}{B4B4B4}
% Demanda externa y comercio, gama fría (azul claro, que se separa del morado)
\definecolor{dicexterna}{HTML}{6AA6DA}
% Zonas: zona del euro (azul marino), España (rojo), resto de países (grises)
\definecolor{dicze}{HTML}{17406D}
\definecolor{dicespana}{HTML}{C8382E}
\definecolor{dicotro}{HTML}{8A8A8A}
\definecolor{dicotroclaro}{HTML}{B5B5B5}
\definecolor{dicotromasclaro}{HTML}{DADADA}
% Precios: inflación general (negro, continuo) y subyacente (negro,
% discontinuo); componentes: energía (bermellón), alimentos (amarillo) y los
% de la subyacente en azules (bienes industriales claro, servicios oscuro).
\definecolor{dicinflacion}{HTML}{1A1A1A}
\definecolor{dicenergia}{HTML}{D55E00}
\definecolor{dicalimentos}{HTML}{F2D45C}
\definecolor{dicbine}{HTML}{7FB2E0}
\definecolor{dicservicios}{HTML}{1F4E89}
% Trabajo: empleo (verde azulado), horas (el mismo, claro), tasa de actividad
% (verde azulado oscuro), paro (teja); productividad (dorado de la web).
\definecolor{dicempleo}{HTML}{1B7F6B}
\definecolor{dichoras}{HTML}{86CBB8}
\definecolor{dicactividad}{HTML}{0F5247}
\definecolor{dicparo}{HTML}{A23B2A}
\definecolor{dicproductividad}{HTML}{B8860B}
% Hogares: renta disponible (pizarra) y tasa de ahorro (negro); el consumo es
% dicconspriv.
\definecolor{dicrenta}{HTML}{4A5A6A}
\definecolor{dicahorro}{HTML}{1A1A1A}
% Política monetaria: tipos del BCE en azules (depósito, el de la zona del
% euro; operaciones principales, medio; facilidad marginal, claro); balance:
% préstamos a los bancos (dorado claro), resto del activo (gris claro).
\definecolor{dicdeposito}{HTML}{17406D}
\definecolor{dicopf}{HTML}{4F86C6}
\definecolor{dicfmc}{HTML}{9EC5EA}
\definecolor{dicprestamos}{HTML}{D9A520}
\definecolor{dicresto}{HTML}{CFCFCF}
% Ciclo: series adelantadas, coincidentes (el PIB) y retrasadas
\definecolor{dicadelantada}{HTML}{D55E00}
\definecolor{dicretrasada}{HTML}{6AA6DA}
% Referencias: rejilla, previsión, sombreados
\definecolor{dicreferencia}{HTML}{767171}

\tikzset{
  % barra de un componente: \addplot[dicbarra=dicfbcf]
  dicbarra/.style={fill=#1, draw=#1!65!black, line width=0.3pt},
  % existencias: gris con trama, para que se distinga también en gris
  dicbarratrama/.style={fill=#1, draw=#1!65!black, line width=0.3pt,
    postaction={pattern=north east lines, pattern color=white}},
  % línea de una serie; con halo blanco cuando pasa sobre barras
  diclinea/.style={very thick, #1},
  diclineahalo/.style={very thick, #1, preaction={draw=white, line width=3.4pt, solid}},
  % puntos de una serie (capa aparte, después de la línea)
  dicpuntos/.style={only marks, mark=*, mark size=1.8pt,
    mark options={fill=white, draw=#1, line width=0.9pt, solid}},
  % línea del cero o de un umbral, después de las barras
  diccero/.style={thin, black},
  dicumbral/.style={thin, dashed, dicreferencia},
}
% Capas de los gráficos de datos (CRITERIOS §7, «tres reglas»), de abajo
% arriba: cuadrícula → sombreados (dicfondo) → barras y áreas (dicbarras) → ejes, cero y umbrales
% (axis lines) → líneas (main) → puntos (dicpuntos) → rótulos (axis
% descriptions). No basta el orden del código: pgfplots dibuja cuadrícula,
% ejes y marcas en sus propias capas (y area style pone los ejes encima), y
% dentro de un axis los \draw, \fill y \node no admiten capa (van con los
% datos de main); por eso todo lo que se dibuja dentro del axis es un \addplot.
% Uso: \begin{axis}[…, dicapas] (al final de las opciones);
%   barras y áreas: \addplot[on layer=dicbarras, …];
%   sombreados:     \dicsombra[estilo]{x1}{y1}{x2}{y2} (rectángulo, en dicfondo,
%                   bajo las barras y sobre la cuadrícula);
%   cero y umbrales: \dicraya[estilo]{(x1,y1) (x2,y2)} (en axis lines);
%   líneas: \addplot normal (main); puntos: \addplot[on layer=dicpuntos, …]
%   (las marcas de cualquier serie van solas a dicpuntos);
%   rótulos: \dicrotulo[opciones del nodo]{(x,y)}{texto} (en axis descriptions),
%   siempre sin fondo (nada de fill=white: si chocan con una línea, se desplazan;
%   CRITERIOS §7, regla 4).
% Los \capa{n} del vídeo no cambian: siguen envolviendo instrucciones completas.
\pgfdeclarelayer{dicfondo}
\pgfdeclarelayer{dicbarras}
\pgfdeclarelayer{dicpuntos}
\pgfplotsset{
  % (el segundo argumento reparte cuadrícula, ejes y rótulos como el estándar)
  layers/dic/.define layer set={axis background, axis grid, dicfondo, dicbarras,
    axis ticks, axis lines, axis tick labels, pre main, main, dicpuntos,
    axis descriptions, axis foreground}{/pgfplots/layers/standard},
  % leyenda sin caja blanca: toma el fondo del gráfico, también si el gráfico va
  % dentro de un cuadro gris (CRITERIOS §7, regla 4, 26/10/2026)
  dicapas/.style={axis on top=false, set layers=dic, mark layer=dicpuntos,
    legend style={fill=none}},
  dicsuelto/.style={forget plot, stack plots=false, sharp plot, mark=none},
}
\newcommand{\dicraya}[2][]{%
  \addplot[on layer=axis lines, dicsuelto, diccero, #1] coordinates {#2};}
\newcommand{\dicsombra}[5][]{%
  \addplot[on layer=dicfondo, dicsuelto, draw=none, #1]
    coordinates {(#2,#3) (#4,#3) (#4,#5) (#2,#5)};}
\newcommand{\dicrotulo}[3][]{%
  \addplot[on layer=axis descriptions, dicsuelto, draw=none] coordinates {#2}
    node[pos=0, #1] {#3};}

%   \begin{document}
%   \begin{tikzpicture}[tcee]
%     \ejes{Q}{P}{6}{5}
%     \draw[curva1] (0.5,4.5) -- (5,0.8) node[etiqueta,right] {$D$};
%     \capa{2}{\draw[curva2] ...;}   % en el vídeo aparece en el paso 2
%   \end{tikzpicture}
%   \end{document}
%
% Paleta: la de la web (morado, dorado, salvia) más tres colores de apoyo
% distinguibles en blanco y negro por el trazo.
% ==========================================================================
\usepackage[T1]{fontenc}
\usepackage{newpxtext,newpxmath}
\usepackage{xcolor}
\usepackage{tikz,pgfplots}
\pgfplotsset{compat=1.18}
\usetikzlibrary{arrows.meta,calc,positioning,patterns,decorations.pathreplacing,intersections,shapes.misc,backgrounds}

\definecolor{tceemorado}{HTML}{5F2987}
\definecolor{tceedorado}{HTML}{B8860B}
\definecolor{tceeverde}{HTML}{3C7D3A}
\definecolor{tceeazul}{HTML}{1F5F99}
\definecolor{tceerojo}{HTML}{B22222}
\definecolor{tceenaranja}{HTML}{E07B28}
\definecolor{tceegris}{HTML}{767171}
\definecolor{tceesalvia}{HTML}{E2EFD9}

% Capas para el vídeo: \capa{n}{…} dibuja su contenido solo a partir del paso n.
% En el PDF, la web y Word se ve todo; generar-video.py compila el gráfico una
% vez por paso con \def\tceepaso{k} para que cada elemento aparezca al narrarlo.
% Vectores con flecha encima, como en el texto (tcee.sty): \vect{x}, \vect{p}
\providecommand{\vect}[1]{\vec{#1}}
\providecommand{\tceepaso}{1000}
\newcommand{\capa}[2]{\ifnum#1>\tceepaso\relax\else#2\fi}

\tikzset{
  tcee/.style={line cap=round,line join=round,font=\small,>={Stealth[length=2.2mm]}},
  eje/.style={->,thick,black},
  curva1/.style={very thick,tceemorado},
  curva2/.style={very thick,tceeazul},
  curva3/.style={very thick,tceeverde},
  curva4/.style={very thick,tceedorado},
  curvanueva/.style={very thick,tceerojo},
  desplazada/.style={very thick,dashed},
  guia/.style={thin,dashed,tceegris},
  punto/.style={circle,fill=black,inner sep=1.3pt},
  etiqueta/.style={font=\small,inner sep=2pt},
  flecha/.style={->,thick,tceerojo},
  area/.style={fill=tceemorado,fill opacity=0.18,draw=none},
  area2/.style={fill=tceedorado,fill opacity=0.22,draw=none},
}

% \ejes{etiqueta x}{etiqueta y}{largo x}{alto y}
\newcommand{\ejes}[4]{%
  \draw[eje] (0,0) -- (#3,0) node[below left=2pt and 0pt] {#1};
  \draw[eje] (0,0) -- (0,#4) node[left] {#2};
  \node[below left] at (0,0) {$0$};}
% \proyeccion{(x,y)}{etq x}{etq y}: líneas guía a los ejes con etiquetas
\newcommand{\proyeccion}[3]{%
  \draw[guia] let \p1=#1 in (\x1,0) node[below] {#2} |- (0,\y1) node[left] {#3};}

% --------------------------------------------------------------------------
% Paleta «dictamen» (dic…): gráficos de datos de la Guía del dictamen (1D01) y
% de cualquier gráfico de coyuntura que quiera los mismos colores. Un concepto,
% un color en todos los gráficos (CRITERIOS §7, «tres reglas», 24/10/2026);
% tabla completa en fuentes/1D01/_trabajo/paleta-1D01.md. No cambia nada de
% lo anterior: son colores y estilos nuevos.
% Actividad: el PIB (y sus variantes: PIB por persona, PIB mundial), morado.
\definecolor{dicpib}{HTML}{5F2987}
% Demanda interna, gama cálida: consumo (naranjas: privado oscuro, público
% claro), inversión (verde) y existencias (gris con trama, aparte).
\definecolor{dicconspriv}{HTML}{C76A1F}
\definecolor{dicconspub}{HTML}{F5C690}
\definecolor{dicfbcf}{HTML}{2A6B3C}
\definecolor{dicexistencias}{HTML}{B4B4B4}
% Demanda externa y comercio, gama fría (azul claro, que se separa del morado)
\definecolor{dicexterna}{HTML}{6AA6DA}
% Zonas: zona del euro (azul marino), España (rojo), resto de países (grises)
\definecolor{dicze}{HTML}{17406D}
\definecolor{dicespana}{HTML}{C8382E}
\definecolor{dicotro}{HTML}{8A8A8A}
\definecolor{dicotroclaro}{HTML}{B5B5B5}
\definecolor{dicotromasclaro}{HTML}{DADADA}
% Precios: inflación general (negro, continuo) y subyacente (negro,
% discontinuo); componentes: energía (bermellón), alimentos (amarillo) y los
% de la subyacente en azules (bienes industriales claro, servicios oscuro).
\definecolor{dicinflacion}{HTML}{1A1A1A}
\definecolor{dicenergia}{HTML}{D55E00}
\definecolor{dicalimentos}{HTML}{F2D45C}
\definecolor{dicbine}{HTML}{7FB2E0}
\definecolor{dicservicios}{HTML}{1F4E89}
% Trabajo: empleo (verde azulado), horas (el mismo, claro), tasa de actividad
% (verde azulado oscuro), paro (teja); productividad (dorado de la web).
\definecolor{dicempleo}{HTML}{1B7F6B}
\definecolor{dichoras}{HTML}{86CBB8}
\definecolor{dicactividad}{HTML}{0F5247}
\definecolor{dicparo}{HTML}{A23B2A}
\definecolor{dicproductividad}{HTML}{B8860B}
% Hogares: renta disponible (pizarra) y tasa de ahorro (negro); el consumo es
% dicconspriv.
\definecolor{dicrenta}{HTML}{4A5A6A}
\definecolor{dicahorro}{HTML}{1A1A1A}
% Política monetaria: tipos del BCE en azules (depósito, el de la zona del
% euro; operaciones principales, medio; facilidad marginal, claro); balance:
% préstamos a los bancos (dorado claro), resto del activo (gris claro).
\definecolor{dicdeposito}{HTML}{17406D}
\definecolor{dicopf}{HTML}{4F86C6}
\definecolor{dicfmc}{HTML}{9EC5EA}
\definecolor{dicprestamos}{HTML}{D9A520}
\definecolor{dicresto}{HTML}{CFCFCF}
% Ciclo: series adelantadas, coincidentes (el PIB) y retrasadas
\definecolor{dicadelantada}{HTML}{D55E00}
\definecolor{dicretrasada}{HTML}{6AA6DA}
% Referencias: rejilla, previsión, sombreados
\definecolor{dicreferencia}{HTML}{767171}

\tikzset{
  % barra de un componente: \addplot[dicbarra=dicfbcf]
  dicbarra/.style={fill=#1, draw=#1!65!black, line width=0.3pt},
  % existencias: gris con trama, para que se distinga también en gris
  dicbarratrama/.style={fill=#1, draw=#1!65!black, line width=0.3pt,
    postaction={pattern=north east lines, pattern color=white}},
  % línea de una serie; con halo blanco cuando pasa sobre barras
  diclinea/.style={very thick, #1},
  diclineahalo/.style={very thick, #1, preaction={draw=white, line width=3.4pt, solid}},
  % puntos de una serie (capa aparte, después de la línea)
  dicpuntos/.style={only marks, mark=*, mark size=1.8pt,
    mark options={fill=white, draw=#1, line width=0.9pt, solid}},
  % línea del cero o de un umbral, después de las barras
  diccero/.style={thin, black},
  dicumbral/.style={thin, dashed, dicreferencia},
}
% Capas de los gráficos de datos (CRITERIOS §7, «tres reglas»), de abajo
% arriba: cuadrícula → sombreados (dicfondo) → barras y áreas (dicbarras) → ejes, cero y umbrales
% (axis lines) → líneas (main) → puntos (dicpuntos) → rótulos (axis
% descriptions). No basta el orden del código: pgfplots dibuja cuadrícula,
% ejes y marcas en sus propias capas (y area style pone los ejes encima), y
% dentro de un axis los \draw, \fill y \node no admiten capa (van con los
% datos de main); por eso todo lo que se dibuja dentro del axis es un \addplot.
% Uso: \begin{axis}[…, dicapas] (al final de las opciones);
%   barras y áreas: \addplot[on layer=dicbarras, …];
%   sombreados:     \dicsombra[estilo]{x1}{y1}{x2}{y2} (rectángulo, en dicfondo,
%                   bajo las barras y sobre la cuadrícula);
%   cero y umbrales: \dicraya[estilo]{(x1,y1) (x2,y2)} (en axis lines);
%   líneas: \addplot normal (main); puntos: \addplot[on layer=dicpuntos, …]
%   (las marcas de cualquier serie van solas a dicpuntos);
%   rótulos: \dicrotulo[opciones del nodo]{(x,y)}{texto} (en axis descriptions),
%   siempre sin fondo (nada de fill=white: si chocan con una línea, se desplazan;
%   CRITERIOS §7, regla 4).
% Los \capa{n} del vídeo no cambian: siguen envolviendo instrucciones completas.
\pgfdeclarelayer{dicfondo}
\pgfdeclarelayer{dicbarras}
\pgfdeclarelayer{dicpuntos}
\pgfplotsset{
  % (el segundo argumento reparte cuadrícula, ejes y rótulos como el estándar)
  layers/dic/.define layer set={axis background, axis grid, dicfondo, dicbarras,
    axis ticks, axis lines, axis tick labels, pre main, main, dicpuntos,
    axis descriptions, axis foreground}{/pgfplots/layers/standard},
  % leyenda sin caja blanca: toma el fondo del gráfico, también si el gráfico va
  % dentro de un cuadro gris (CRITERIOS §7, regla 4, 26/10/2026)
  dicapas/.style={axis on top=false, set layers=dic, mark layer=dicpuntos,
    legend style={fill=none}},
  dicsuelto/.style={forget plot, stack plots=false, sharp plot, mark=none},
}
\newcommand{\dicraya}[2][]{%
  \addplot[on layer=axis lines, dicsuelto, diccero, #1] coordinates {#2};}
\newcommand{\dicsombra}[5][]{%
  \addplot[on layer=dicfondo, dicsuelto, draw=none, #1]
    coordinates {(#2,#3) (#4,#3) (#4,#5) (#2,#5)};}
\newcommand{\dicrotulo}[3][]{%
  \addplot[on layer=axis descriptions, dicsuelto, draw=none] coordinates {#2}
    node[pos=0, #1] {#3};}

% Variantes: \def\cajacapas{1,2,6} antes de % Caja de Edgeworth del intercambio puro (3.A.21): gráfico maestro que crece por
% capas desde I.3 hasta III.2. Las variantes del PDF (caja-intercambio-equilibrio,
% caja-intercambio-teoremas…) fijan \cajacapas con las capas que muestran y
% cargan este fichero: mismo código y mismas cifras.
%
% Dotaciones totales X̄ = 10, Ȳ = 6; dotación ω = (1,5; 5) para A y (8,5; 1) para B.
% U_A = X_A^0,4 Y_A^0,6 (desde O_A); U_B = X_B^0,6 Y_B^0,4 (desde O_B).
% RMS^A_XY = (2/3) Y_A/X_A; RMS^B_XY = 1,5 Y_B/X_B. En ω: 2,22 y 0,18 (lente).
% Demandas (p_Y = 1, p = p_X): X_A = 0,4 (1,5p + 5)/p, Y_A = 0,6 (1,5p + 5);
%   X_B = 0,6 (8,5p + 1)/p, Y_B = 0,4 (8,5p + 1) (en la caja, X̄ - X_B, Ȳ - Y_B).
% Curvas de oferta-demanda: esos puntos al variar p.
% Equilibrio: p* = (0,4·5 + 0,6·1)/(0,6·1,5 + 0,4·8,5) = 0,6047;
%   E = (3,908; 3,544); U_A(E) = 3,685; U_B(E) = 4,236.
% Desequilibrio p = 0,23: A elige (9,296; 3,207) y B (2,291; 4,818) en la caja:
%   exceso de demanda de X = 7,004 y exceso de oferta de Y = 1,611 (0,23·7,004).
% Curva de contrato: Y_A = 9 X_A / ((2/3)(10 - X_A) + 1,5 X_A).
% Lente: U_A = U_A(ω) = 3,090 y U_B = U_B(ω) = 3,610; se cortan en X_A = 1,5 y 6,713.
% CC' (curva de contrato dentro de la lente): C = (3,148; 3,050), C' = (4,675; 3,983).
% 2TFEB: E' = (2,2; 2,329) en la curva de contrato, RMS = 0,706; la recta de
%   precios tangente en E' pasa por ω' = (1,5; 2,824): A cede 2,18 de Y a B.
% Autor: Víctor Gutiérrez Marcos
\documentclass[tikz,border=4pt]{standalone}
% ==========================================================================
% tikz-tcee.tex — Estilos comunes de los gráficos TikZ del temario
% Autor: Víctor Gutiérrez Marcos
%
% Cada gráfico es un fichero standalone en fuentes/<tema>/graficos/:
%
%   \documentclass[tikz,border=4pt]{standalone}
%   \input{../../../latex/tikz-tcee}
%   \begin{document}
%   \begin{tikzpicture}[tcee]
%     \ejes{Q}{P}{6}{5}
%     \draw[curva1] (0.5,4.5) -- (5,0.8) node[etiqueta,right] {$D$};
%     \capa{2}{\draw[curva2] ...;}   % en el vídeo aparece en el paso 2
%   \end{tikzpicture}
%   \end{document}
%
% Paleta: la de la web (morado, dorado, salvia) más tres colores de apoyo
% distinguibles en blanco y negro por el trazo.
% ==========================================================================
\usepackage[T1]{fontenc}
\usepackage{newpxtext,newpxmath}
\usepackage{xcolor}
\usepackage{tikz,pgfplots}
\pgfplotsset{compat=1.18}
\usetikzlibrary{arrows.meta,calc,positioning,patterns,decorations.pathreplacing,intersections,shapes.misc,backgrounds}

\definecolor{tceemorado}{HTML}{5F2987}
\definecolor{tceedorado}{HTML}{B8860B}
\definecolor{tceeverde}{HTML}{3C7D3A}
\definecolor{tceeazul}{HTML}{1F5F99}
\definecolor{tceerojo}{HTML}{B22222}
\definecolor{tceenaranja}{HTML}{E07B28}
\definecolor{tceegris}{HTML}{767171}
\definecolor{tceesalvia}{HTML}{E2EFD9}

% Capas para el vídeo: \capa{n}{…} dibuja su contenido solo a partir del paso n.
% En el PDF, la web y Word se ve todo; generar-video.py compila el gráfico una
% vez por paso con \def\tceepaso{k} para que cada elemento aparezca al narrarlo.
% Vectores con flecha encima, como en el texto (tcee.sty): \vect{x}, \vect{p}
\providecommand{\vect}[1]{\vec{#1}}
\providecommand{\tceepaso}{1000}
\newcommand{\capa}[2]{\ifnum#1>\tceepaso\relax\else#2\fi}

\tikzset{
  tcee/.style={line cap=round,line join=round,font=\small,>={Stealth[length=2.2mm]}},
  eje/.style={->,thick,black},
  curva1/.style={very thick,tceemorado},
  curva2/.style={very thick,tceeazul},
  curva3/.style={very thick,tceeverde},
  curva4/.style={very thick,tceedorado},
  curvanueva/.style={very thick,tceerojo},
  desplazada/.style={very thick,dashed},
  guia/.style={thin,dashed,tceegris},
  punto/.style={circle,fill=black,inner sep=1.3pt},
  etiqueta/.style={font=\small,inner sep=2pt},
  flecha/.style={->,thick,tceerojo},
  area/.style={fill=tceemorado,fill opacity=0.18,draw=none},
  area2/.style={fill=tceedorado,fill opacity=0.22,draw=none},
}

% \ejes{etiqueta x}{etiqueta y}{largo x}{alto y}
\newcommand{\ejes}[4]{%
  \draw[eje] (0,0) -- (#3,0) node[below left=2pt and 0pt] {#1};
  \draw[eje] (0,0) -- (0,#4) node[left] {#2};
  \node[below left] at (0,0) {$0$};}
% \proyeccion{(x,y)}{etq x}{etq y}: líneas guía a los ejes con etiquetas
\newcommand{\proyeccion}[3]{%
  \draw[guia] let \p1=#1 in (\x1,0) node[below] {#2} |- (0,\y1) node[left] {#3};}

% --------------------------------------------------------------------------
% Paleta «dictamen» (dic…): gráficos de datos de la Guía del dictamen (1D01) y
% de cualquier gráfico de coyuntura que quiera los mismos colores. Un concepto,
% un color en todos los gráficos (CRITERIOS §7, «tres reglas», 24/10/2026);
% tabla completa en fuentes/1D01/_trabajo/paleta-1D01.md. No cambia nada de
% lo anterior: son colores y estilos nuevos.
% Actividad: el PIB (y sus variantes: PIB por persona, PIB mundial), morado.
\definecolor{dicpib}{HTML}{5F2987}
% Demanda interna, gama cálida: consumo (naranjas: privado oscuro, público
% claro), inversión (verde) y existencias (gris con trama, aparte).
\definecolor{dicconspriv}{HTML}{C76A1F}
\definecolor{dicconspub}{HTML}{F5C690}
\definecolor{dicfbcf}{HTML}{2A6B3C}
\definecolor{dicexistencias}{HTML}{B4B4B4}
% Demanda externa y comercio, gama fría (azul claro, que se separa del morado)
\definecolor{dicexterna}{HTML}{6AA6DA}
% Zonas: zona del euro (azul marino), España (rojo), resto de países (grises)
\definecolor{dicze}{HTML}{17406D}
\definecolor{dicespana}{HTML}{C8382E}
\definecolor{dicotro}{HTML}{8A8A8A}
\definecolor{dicotroclaro}{HTML}{B5B5B5}
\definecolor{dicotromasclaro}{HTML}{DADADA}
% Precios: inflación general (negro, continuo) y subyacente (negro,
% discontinuo); componentes: energía (bermellón), alimentos (amarillo) y los
% de la subyacente en azules (bienes industriales claro, servicios oscuro).
\definecolor{dicinflacion}{HTML}{1A1A1A}
\definecolor{dicenergia}{HTML}{D55E00}
\definecolor{dicalimentos}{HTML}{F2D45C}
\definecolor{dicbine}{HTML}{7FB2E0}
\definecolor{dicservicios}{HTML}{1F4E89}
% Trabajo: empleo (verde azulado), horas (el mismo, claro), tasa de actividad
% (verde azulado oscuro), paro (teja); productividad (dorado de la web).
\definecolor{dicempleo}{HTML}{1B7F6B}
\definecolor{dichoras}{HTML}{86CBB8}
\definecolor{dicactividad}{HTML}{0F5247}
\definecolor{dicparo}{HTML}{A23B2A}
\definecolor{dicproductividad}{HTML}{B8860B}
% Hogares: renta disponible (pizarra) y tasa de ahorro (negro); el consumo es
% dicconspriv.
\definecolor{dicrenta}{HTML}{4A5A6A}
\definecolor{dicahorro}{HTML}{1A1A1A}
% Política monetaria: tipos del BCE en azules (depósito, el de la zona del
% euro; operaciones principales, medio; facilidad marginal, claro); balance:
% préstamos a los bancos (dorado claro), resto del activo (gris claro).
\definecolor{dicdeposito}{HTML}{17406D}
\definecolor{dicopf}{HTML}{4F86C6}
\definecolor{dicfmc}{HTML}{9EC5EA}
\definecolor{dicprestamos}{HTML}{D9A520}
\definecolor{dicresto}{HTML}{CFCFCF}
% Ciclo: series adelantadas, coincidentes (el PIB) y retrasadas
\definecolor{dicadelantada}{HTML}{D55E00}
\definecolor{dicretrasada}{HTML}{6AA6DA}
% Referencias: rejilla, previsión, sombreados
\definecolor{dicreferencia}{HTML}{767171}

\tikzset{
  % barra de un componente: \addplot[dicbarra=dicfbcf]
  dicbarra/.style={fill=#1, draw=#1!65!black, line width=0.3pt},
  % existencias: gris con trama, para que se distinga también en gris
  dicbarratrama/.style={fill=#1, draw=#1!65!black, line width=0.3pt,
    postaction={pattern=north east lines, pattern color=white}},
  % línea de una serie; con halo blanco cuando pasa sobre barras
  diclinea/.style={very thick, #1},
  diclineahalo/.style={very thick, #1, preaction={draw=white, line width=3.4pt, solid}},
  % puntos de una serie (capa aparte, después de la línea)
  dicpuntos/.style={only marks, mark=*, mark size=1.8pt,
    mark options={fill=white, draw=#1, line width=0.9pt, solid}},
  % línea del cero o de un umbral, después de las barras
  diccero/.style={thin, black},
  dicumbral/.style={thin, dashed, dicreferencia},
}
% Capas de los gráficos de datos (CRITERIOS §7, «tres reglas»), de abajo
% arriba: cuadrícula → sombreados (dicfondo) → barras y áreas (dicbarras) → ejes, cero y umbrales
% (axis lines) → líneas (main) → puntos (dicpuntos) → rótulos (axis
% descriptions). No basta el orden del código: pgfplots dibuja cuadrícula,
% ejes y marcas en sus propias capas (y area style pone los ejes encima), y
% dentro de un axis los \draw, \fill y \node no admiten capa (van con los
% datos de main); por eso todo lo que se dibuja dentro del axis es un \addplot.
% Uso: \begin{axis}[…, dicapas] (al final de las opciones);
%   barras y áreas: \addplot[on layer=dicbarras, …];
%   sombreados:     \dicsombra[estilo]{x1}{y1}{x2}{y2} (rectángulo, en dicfondo,
%                   bajo las barras y sobre la cuadrícula);
%   cero y umbrales: \dicraya[estilo]{(x1,y1) (x2,y2)} (en axis lines);
%   líneas: \addplot normal (main); puntos: \addplot[on layer=dicpuntos, …]
%   (las marcas de cualquier serie van solas a dicpuntos);
%   rótulos: \dicrotulo[opciones del nodo]{(x,y)}{texto} (en axis descriptions),
%   siempre sin fondo (nada de fill=white: si chocan con una línea, se desplazan;
%   CRITERIOS §7, regla 4).
% Los \capa{n} del vídeo no cambian: siguen envolviendo instrucciones completas.
\pgfdeclarelayer{dicfondo}
\pgfdeclarelayer{dicbarras}
\pgfdeclarelayer{dicpuntos}
\pgfplotsset{
  % (el segundo argumento reparte cuadrícula, ejes y rótulos como el estándar)
  layers/dic/.define layer set={axis background, axis grid, dicfondo, dicbarras,
    axis ticks, axis lines, axis tick labels, pre main, main, dicpuntos,
    axis descriptions, axis foreground}{/pgfplots/layers/standard},
  % leyenda sin caja blanca: toma el fondo del gráfico, también si el gráfico va
  % dentro de un cuadro gris (CRITERIOS §7, regla 4, 26/10/2026)
  dicapas/.style={axis on top=false, set layers=dic, mark layer=dicpuntos,
    legend style={fill=none}},
  dicsuelto/.style={forget plot, stack plots=false, sharp plot, mark=none},
}
\newcommand{\dicraya}[2][]{%
  \addplot[on layer=axis lines, dicsuelto, diccero, #1] coordinates {#2};}
\newcommand{\dicsombra}[5][]{%
  \addplot[on layer=dicfondo, dicsuelto, draw=none, #1]
    coordinates {(#2,#3) (#4,#3) (#4,#5) (#2,#5)};}
\newcommand{\dicrotulo}[3][]{%
  \addplot[on layer=axis descriptions, dicsuelto, draw=none] coordinates {#2}
    node[pos=0, #1] {#3};}

% Variantes: \def\cajacapas{1,2,6} antes de % Caja de Edgeworth del intercambio puro (3.A.21): gráfico maestro que crece por
% capas desde I.3 hasta III.2. Las variantes del PDF (caja-intercambio-equilibrio,
% caja-intercambio-teoremas…) fijan \cajacapas con las capas que muestran y
% cargan este fichero: mismo código y mismas cifras.
%
% Dotaciones totales X̄ = 10, Ȳ = 6; dotación ω = (1,5; 5) para A y (8,5; 1) para B.
% U_A = X_A^0,4 Y_A^0,6 (desde O_A); U_B = X_B^0,6 Y_B^0,4 (desde O_B).
% RMS^A_XY = (2/3) Y_A/X_A; RMS^B_XY = 1,5 Y_B/X_B. En ω: 2,22 y 0,18 (lente).
% Demandas (p_Y = 1, p = p_X): X_A = 0,4 (1,5p + 5)/p, Y_A = 0,6 (1,5p + 5);
%   X_B = 0,6 (8,5p + 1)/p, Y_B = 0,4 (8,5p + 1) (en la caja, X̄ - X_B, Ȳ - Y_B).
% Curvas de oferta-demanda: esos puntos al variar p.
% Equilibrio: p* = (0,4·5 + 0,6·1)/(0,6·1,5 + 0,4·8,5) = 0,6047;
%   E = (3,908; 3,544); U_A(E) = 3,685; U_B(E) = 4,236.
% Desequilibrio p = 0,23: A elige (9,296; 3,207) y B (2,291; 4,818) en la caja:
%   exceso de demanda de X = 7,004 y exceso de oferta de Y = 1,611 (0,23·7,004).
% Curva de contrato: Y_A = 9 X_A / ((2/3)(10 - X_A) + 1,5 X_A).
% Lente: U_A = U_A(ω) = 3,090 y U_B = U_B(ω) = 3,610; se cortan en X_A = 1,5 y 6,713.
% CC' (curva de contrato dentro de la lente): C = (3,148; 3,050), C' = (4,675; 3,983).
% 2TFEB: E' = (2,2; 2,329) en la curva de contrato, RMS = 0,706; la recta de
%   precios tangente en E' pasa por ω' = (1,5; 2,824): A cede 2,18 de Y a B.
% Autor: Víctor Gutiérrez Marcos
\documentclass[tikz,border=4pt]{standalone}
\input{../../../latex/tikz-tcee}
% Variantes: \def\cajacapas{1,2,6} antes de \input{caja-intercambio} muestra solo
% esas capas (en el vídeo, \tceepaso sigue mandando).
\makeatletter
\ifdefined\cajacapas
  \let\capamaestro\capa
  \renewcommand{\capa}[2]{\edef\cajaesta{\noexpand\in@{,#1,}{,\cajacapas,}}\cajaesta
    \ifin@\capamaestro{#1}{#2}\fi}
\fi
\makeatother
\begin{document}
\begin{tikzpicture}[tcee,x=1cm,y=1.1cm,
    declare function={
      ua(\x,\u)=(\u/\x^0.4)^(1/0.6);
      ub(\x,\v)=6-(\v/(10-\x)^0.6)^(1/0.4);
      cc(\x)=9*\x/((2/3)*(10-\x)+1.5*\x);
      odax(\p)=0.4*(1.5*\p+5)/\p;  oday(\p)=0.6*(1.5*\p+5);
      odbx(\p)=10-0.6*(8.5*\p+1)/\p; odby(\p)=6-0.4*(8.5*\p+1);},
    llave/.style={decorate,decoration={brace,amplitude=5pt},thin,tceegris}]
  % 1: la caja, sus orígenes, los ejes de cada consumidor y la dotación
  \capa{1}{
    \draw[thick] (0,0) rectangle (10,6);
    \node[below left,text=tceemorado] at (0,0) {$O_A$};
    \node[above right,text=tceeverde] at (10,6) {$O_B$};
    \draw[->,thin,tceemorado] (0.4,-0.3) -- (1.4,-0.3) node[right,etiqueta] {$X_A$};
    \draw[->,thin,tceemorado] (-0.3,0.4) -- (-0.3,1.4) node[above,etiqueta] {$Y_A$};
    \draw[->,thin,tceeverde] (9.6,6.3) -- (8.6,6.3) node[left,etiqueta] {$X_B$};
    \draw[->,thin,tceeverde] (10.3,5.6) -- (10.3,4.6) node[below,etiqueta] {$Y_B$};
    \draw[llave] (10,-0.75) -- (0,-0.75) node[midway,below=5pt,text=black] {$\bar X$};
    \draw[llave] (-0.95,0) -- (-0.95,6) node[midway,left=5pt,text=black] {$\bar Y$};
  }
  \capa{1}{\node[punto,label={[etiqueta]below left:$\vect{\omega}$}] (w) at (1.5,5) {};}
  % 2: curvas de indiferencia que pasan por ω y la lente
  \capa{2}{\draw[curva1] plot[domain=1.2:9.3,samples=120] (\x,{ua(\x,3.09017)})
    node[etiqueta,right] {$U_A^0$};}
  \capa{2}{\draw[curva3] plot[domain=0.2:7.38,samples=120] (\x,{ub(\x,3.60986)})
    node[etiqueta,right] {$U_B^0$};}
  \capa{2}{\fill[area,fill opacity=0.12] plot[domain=1.5:6.7133,samples=100] (\x,{ua(\x,3.09017)})
    -- plot[domain=6.7133:1.5,samples=100] (\x,{ub(\x,3.60986)}) -- cycle;}
  % 3: precios de desequilibrio (p_X/p_Y = 0,23): lo que elige cada uno y los excesos
  \capa{3}{\draw[curva4,desplazada] (0.5,5.23) -- (9.6,3.137) node[etiqueta,anchor=north east] {$p_X/p_Y$};}
  \capa{3}{\node[punto,label={[etiqueta]above left:$\vect{x}_A$}] at (9.2957,3.207) {};}
  \capa{3}{\node[punto,label={[etiqueta,label distance=3pt]above:$\vect{x}_B$}] at (2.2913,4.818) {};}
  \capa{3}{\draw[curvanueva] (2.2913,4.818) -- (9.2957,4.818) node[midway,above,etiqueta] {$z_X>0$};}
  \capa{3}{\draw[curvanueva] (9.2957,4.818) -- (9.2957,3.207) node[pos=0.3,left,etiqueta] {$z_Y<0$};}
  % 4: curva de oferta-demanda de A
  \capa{4}{\draw[curva1,desplazada,thick] plot[domain=2.2222:0.23,samples=150,variable=\p]
    ({odax(\p)},{oday(\p)});}
  \capa{4}{\node[etiqueta,text=tceemorado,above] at (6.66,3.297) {$\mathrm{OD}_A$};}
  % 5: curva de oferta-demanda de B
  \capa{5}{\draw[curva3,desplazada,thick] plot[domain=0.17647:1.25,samples=150,variable=\p]
    ({odbx(\p)},{odby(\p)}) node[etiqueta,below] {$\mathrm{OD}_B$};}
  % 6: equilibrio E, recta de precios de equilibrio y curvas tangentes en E
  \capa{6}{\node[punto,label={[etiqueta]above:$E$}] (E) at (3.9077,3.5442) {};}
  \capa{6}{\draw[curva4] (0.6,5.5442) -- (9.7686,0) node[etiqueta,below] {$p^*_X/p^*_Y$};}
  \capa{6}{\draw[curva1] plot[domain=1.9:9.3,samples=120] (\x,{ua(\x,3.68534)})
    node[etiqueta,right] {$U_A^*$};}
  \capa{6}{\draw[curva3] plot[domain=0.2:6.6,samples=120] (\x,{ub(\x,4.23592)})
    node[etiqueta,right] {$U_B^*$};}
  % 7: curva de contrato (1TFEB: E está en ella)
  \capa{7}{\draw[curva2] plot[domain=0:10,samples=150] (\x,{cc(\x)});}
  \capa{7}{\node[etiqueta,text=tceeazul,anchor=east] at (7.9,5.75) {curva de contrato};}
  % 8: 2TFEB: otro óptimo E', su recta de precios y la redistribución ω → ω'
  \capa{8}{\node[punto,label={[etiqueta]right:$E'$}] at (2.2,2.3294) {};}
  \capa{8}{\draw[curva4] (0.9,3.2471) -- (5.5,0) ;}
  \capa{8}{\node[punto,label={[etiqueta]left:$\vect{\omega}'$}] (w2) at (1.5,2.8235) {};}
  \capa{8}{\draw[flecha,shorten <=3pt,shorten >=3pt] (w) -- (w2);}
  % 9: Edgeworth: el tramo CC' de la curva de contrato dentro de la lente
  \capa{9}{\draw[curva2,line width=5pt,opacity=0.35] plot[domain=3.1484:4.6748,samples=40] (\x,{cc(\x)});}
  \capa{9}{\node[punto,label={[etiqueta]left:$C$}] at (3.1484,3.0500) {};}
  \capa{9}{\node[punto,label={[etiqueta]right:$C'$}] at (4.6748,3.9834) {};}
  % 10: el núcleo se encoge hacia E al replicar la economía
  \capa{10}{\draw[flecha] plot[domain=3.2:3.78,samples=20] (\x,{cc(\x)});}
  \capa{10}{\draw[flecha] plot[domain=4.6:4.04,samples=20] (\x,{cc(\x)});}
  \capa{10}{\node[etiqueta,text=tceerojo] at (3.5,2.45) {$r\to\infty$};}
\end{tikzpicture}
\end{document}
 muestra solo
% esas capas (en el vídeo, \tceepaso sigue mandando).
\makeatletter
\ifdefined\cajacapas
  \let\capamaestro\capa
  \renewcommand{\capa}[2]{\edef\cajaesta{\noexpand\in@{,#1,}{,\cajacapas,}}\cajaesta
    \ifin@\capamaestro{#1}{#2}\fi}
\fi
\makeatother
\begin{document}
\begin{tikzpicture}[tcee,x=1cm,y=1.1cm,
    declare function={
      ua(\x,\u)=(\u/\x^0.4)^(1/0.6);
      ub(\x,\v)=6-(\v/(10-\x)^0.6)^(1/0.4);
      cc(\x)=9*\x/((2/3)*(10-\x)+1.5*\x);
      odax(\p)=0.4*(1.5*\p+5)/\p;  oday(\p)=0.6*(1.5*\p+5);
      odbx(\p)=10-0.6*(8.5*\p+1)/\p; odby(\p)=6-0.4*(8.5*\p+1);},
    llave/.style={decorate,decoration={brace,amplitude=5pt},thin,tceegris}]
  % 1: la caja, sus orígenes, los ejes de cada consumidor y la dotación
  \capa{1}{
    \draw[thick] (0,0) rectangle (10,6);
    \node[below left,text=tceemorado] at (0,0) {$O_A$};
    \node[above right,text=tceeverde] at (10,6) {$O_B$};
    \draw[->,thin,tceemorado] (0.4,-0.3) -- (1.4,-0.3) node[right,etiqueta] {$X_A$};
    \draw[->,thin,tceemorado] (-0.3,0.4) -- (-0.3,1.4) node[above,etiqueta] {$Y_A$};
    \draw[->,thin,tceeverde] (9.6,6.3) -- (8.6,6.3) node[left,etiqueta] {$X_B$};
    \draw[->,thin,tceeverde] (10.3,5.6) -- (10.3,4.6) node[below,etiqueta] {$Y_B$};
    \draw[llave] (10,-0.75) -- (0,-0.75) node[midway,below=5pt,text=black] {$\bar X$};
    \draw[llave] (-0.95,0) -- (-0.95,6) node[midway,left=5pt,text=black] {$\bar Y$};
  }
  \capa{1}{\node[punto,label={[etiqueta]below left:$\vect{\omega}$}] (w) at (1.5,5) {};}
  % 2: curvas de indiferencia que pasan por ω y la lente
  \capa{2}{\draw[curva1] plot[domain=1.2:9.3,samples=120] (\x,{ua(\x,3.09017)})
    node[etiqueta,right] {$U_A^0$};}
  \capa{2}{\draw[curva3] plot[domain=0.2:7.38,samples=120] (\x,{ub(\x,3.60986)})
    node[etiqueta,right] {$U_B^0$};}
  \capa{2}{\fill[area,fill opacity=0.12] plot[domain=1.5:6.7133,samples=100] (\x,{ua(\x,3.09017)})
    -- plot[domain=6.7133:1.5,samples=100] (\x,{ub(\x,3.60986)}) -- cycle;}
  % 3: precios de desequilibrio (p_X/p_Y = 0,23): lo que elige cada uno y los excesos
  \capa{3}{\draw[curva4,desplazada] (0.5,5.23) -- (9.6,3.137) node[etiqueta,anchor=north east] {$p_X/p_Y$};}
  \capa{3}{\node[punto,label={[etiqueta]above left:$\vect{x}_A$}] at (9.2957,3.207) {};}
  \capa{3}{\node[punto,label={[etiqueta,label distance=3pt]above:$\vect{x}_B$}] at (2.2913,4.818) {};}
  \capa{3}{\draw[curvanueva] (2.2913,4.818) -- (9.2957,4.818) node[midway,above,etiqueta] {$z_X>0$};}
  \capa{3}{\draw[curvanueva] (9.2957,4.818) -- (9.2957,3.207) node[pos=0.3,left,etiqueta] {$z_Y<0$};}
  % 4: curva de oferta-demanda de A
  \capa{4}{\draw[curva1,desplazada,thick] plot[domain=2.2222:0.23,samples=150,variable=\p]
    ({odax(\p)},{oday(\p)});}
  \capa{4}{\node[etiqueta,text=tceemorado,above] at (6.66,3.297) {$\mathrm{OD}_A$};}
  % 5: curva de oferta-demanda de B
  \capa{5}{\draw[curva3,desplazada,thick] plot[domain=0.17647:1.25,samples=150,variable=\p]
    ({odbx(\p)},{odby(\p)}) node[etiqueta,below] {$\mathrm{OD}_B$};}
  % 6: equilibrio E, recta de precios de equilibrio y curvas tangentes en E
  \capa{6}{\node[punto,label={[etiqueta]above:$E$}] (E) at (3.9077,3.5442) {};}
  \capa{6}{\draw[curva4] (0.6,5.5442) -- (9.7686,0) node[etiqueta,below] {$p^*_X/p^*_Y$};}
  \capa{6}{\draw[curva1] plot[domain=1.9:9.3,samples=120] (\x,{ua(\x,3.68534)})
    node[etiqueta,right] {$U_A^*$};}
  \capa{6}{\draw[curva3] plot[domain=0.2:6.6,samples=120] (\x,{ub(\x,4.23592)})
    node[etiqueta,right] {$U_B^*$};}
  % 7: curva de contrato (1TFEB: E está en ella)
  \capa{7}{\draw[curva2] plot[domain=0:10,samples=150] (\x,{cc(\x)});}
  \capa{7}{\node[etiqueta,text=tceeazul,anchor=east] at (7.9,5.75) {curva de contrato};}
  % 8: 2TFEB: otro óptimo E', su recta de precios y la redistribución ω → ω'
  \capa{8}{\node[punto,label={[etiqueta]right:$E'$}] at (2.2,2.3294) {};}
  \capa{8}{\draw[curva4] (0.9,3.2471) -- (5.5,0) ;}
  \capa{8}{\node[punto,label={[etiqueta]left:$\vect{\omega}'$}] (w2) at (1.5,2.8235) {};}
  \capa{8}{\draw[flecha,shorten <=3pt,shorten >=3pt] (w) -- (w2);}
  % 9: Edgeworth: el tramo CC' de la curva de contrato dentro de la lente
  \capa{9}{\draw[curva2,line width=5pt,opacity=0.35] plot[domain=3.1484:4.6748,samples=40] (\x,{cc(\x)});}
  \capa{9}{\node[punto,label={[etiqueta]left:$C$}] at (3.1484,3.0500) {};}
  \capa{9}{\node[punto,label={[etiqueta]right:$C'$}] at (4.6748,3.9834) {};}
  % 10: el núcleo se encoge hacia E al replicar la economía
  \capa{10}{\draw[flecha] plot[domain=3.2:3.78,samples=20] (\x,{cc(\x)});}
  \capa{10}{\draw[flecha] plot[domain=4.6:4.04,samples=20] (\x,{cc(\x)});}
  \capa{10}{\node[etiqueta,text=tceerojo] at (3.5,2.45) {$r\to\infty$};}
\end{tikzpicture}
\end{document}
 muestra solo
% esas capas (en el vídeo, \tceepaso sigue mandando).
\makeatletter
\ifdefined\cajacapas
  \let\capamaestro\capa
  \renewcommand{\capa}[2]{\edef\cajaesta{\noexpand\in@{,#1,}{,\cajacapas,}}\cajaesta
    \ifin@\capamaestro{#1}{#2}\fi}
\fi
\makeatother
\begin{document}
\begin{tikzpicture}[tcee,x=1cm,y=1.1cm,
    declare function={
      ua(\x,\u)=(\u/\x^0.4)^(1/0.6);
      ub(\x,\v)=6-(\v/(10-\x)^0.6)^(1/0.4);
      cc(\x)=9*\x/((2/3)*(10-\x)+1.5*\x);
      odax(\p)=0.4*(1.5*\p+5)/\p;  oday(\p)=0.6*(1.5*\p+5);
      odbx(\p)=10-0.6*(8.5*\p+1)/\p; odby(\p)=6-0.4*(8.5*\p+1);},
    llave/.style={decorate,decoration={brace,amplitude=5pt},thin,tceegris}]
  % 1: la caja, sus orígenes, los ejes de cada consumidor y la dotación
  \capa{1}{
    \draw[thick] (0,0) rectangle (10,6);
    \node[below left,text=tceemorado] at (0,0) {$O_A$};
    \node[above right,text=tceeverde] at (10,6) {$O_B$};
    \draw[->,thin,tceemorado] (0.4,-0.3) -- (1.4,-0.3) node[right,etiqueta] {$X_A$};
    \draw[->,thin,tceemorado] (-0.3,0.4) -- (-0.3,1.4) node[above,etiqueta] {$Y_A$};
    \draw[->,thin,tceeverde] (9.6,6.3) -- (8.6,6.3) node[left,etiqueta] {$X_B$};
    \draw[->,thin,tceeverde] (10.3,5.6) -- (10.3,4.6) node[below,etiqueta] {$Y_B$};
    \draw[llave] (10,-0.75) -- (0,-0.75) node[midway,below=5pt,text=black] {$\bar X$};
    \draw[llave] (-0.95,0) -- (-0.95,6) node[midway,left=5pt,text=black] {$\bar Y$};
  }
  \capa{1}{\node[punto,label={[etiqueta]below left:$\vect{\omega}$}] (w) at (1.5,5) {};}
  % 2: curvas de indiferencia que pasan por ω y la lente
  \capa{2}{\draw[curva1] plot[domain=1.2:9.3,samples=120] (\x,{ua(\x,3.09017)})
    node[etiqueta,right] {$U_A^0$};}
  \capa{2}{\draw[curva3] plot[domain=0.2:7.38,samples=120] (\x,{ub(\x,3.60986)})
    node[etiqueta,right] {$U_B^0$};}
  \capa{2}{\fill[area,fill opacity=0.12] plot[domain=1.5:6.7133,samples=100] (\x,{ua(\x,3.09017)})
    -- plot[domain=6.7133:1.5,samples=100] (\x,{ub(\x,3.60986)}) -- cycle;}
  % 3: precios de desequilibrio (p_X/p_Y = 0,23): lo que elige cada uno y los excesos
  \capa{3}{\draw[curva4,desplazada] (0.5,5.23) -- (9.6,3.137) node[etiqueta,anchor=north east] {$p_X/p_Y$};}
  \capa{3}{\node[punto,label={[etiqueta]above left:$\vect{x}_A$}] at (9.2957,3.207) {};}
  \capa{3}{\node[punto,label={[etiqueta,label distance=3pt]above:$\vect{x}_B$}] at (2.2913,4.818) {};}
  \capa{3}{\draw[curvanueva] (2.2913,4.818) -- (9.2957,4.818) node[midway,above,etiqueta] {$z_X>0$};}
  \capa{3}{\draw[curvanueva] (9.2957,4.818) -- (9.2957,3.207) node[pos=0.3,left,etiqueta] {$z_Y<0$};}
  % 4: curva de oferta-demanda de A
  \capa{4}{\draw[curva1,desplazada,thick] plot[domain=2.2222:0.23,samples=150,variable=\p]
    ({odax(\p)},{oday(\p)});}
  \capa{4}{\node[etiqueta,text=tceemorado,above] at (6.66,3.297) {$\mathrm{OD}_A$};}
  % 5: curva de oferta-demanda de B
  \capa{5}{\draw[curva3,desplazada,thick] plot[domain=0.17647:1.25,samples=150,variable=\p]
    ({odbx(\p)},{odby(\p)}) node[etiqueta,below] {$\mathrm{OD}_B$};}
  % 6: equilibrio E, recta de precios de equilibrio y curvas tangentes en E
  \capa{6}{\node[punto,label={[etiqueta]above:$E$}] (E) at (3.9077,3.5442) {};}
  \capa{6}{\draw[curva4] (0.6,5.5442) -- (9.7686,0) node[etiqueta,below] {$p^*_X/p^*_Y$};}
  \capa{6}{\draw[curva1] plot[domain=1.9:9.3,samples=120] (\x,{ua(\x,3.68534)})
    node[etiqueta,right] {$U_A^*$};}
  \capa{6}{\draw[curva3] plot[domain=0.2:6.6,samples=120] (\x,{ub(\x,4.23592)})
    node[etiqueta,right] {$U_B^*$};}
  % 7: curva de contrato (1TFEB: E está en ella)
  \capa{7}{\draw[curva2] plot[domain=0:10,samples=150] (\x,{cc(\x)});}
  \capa{7}{\node[etiqueta,text=tceeazul,anchor=east] at (7.9,5.75) {curva de contrato};}
  % 8: 2TFEB: otro óptimo E', su recta de precios y la redistribución ω → ω'
  \capa{8}{\node[punto,label={[etiqueta]right:$E'$}] at (2.2,2.3294) {};}
  \capa{8}{\draw[curva4] (0.9,3.2471) -- (5.5,0) ;}
  \capa{8}{\node[punto,label={[etiqueta]left:$\vect{\omega}'$}] (w2) at (1.5,2.8235) {};}
  \capa{8}{\draw[flecha,shorten <=3pt,shorten >=3pt] (w) -- (w2);}
  % 9: Edgeworth: el tramo CC' de la curva de contrato dentro de la lente
  \capa{9}{\draw[curva2,line width=5pt,opacity=0.35] plot[domain=3.1484:4.6748,samples=40] (\x,{cc(\x)});}
  \capa{9}{\node[punto,label={[etiqueta]left:$C$}] at (3.1484,3.0500) {};}
  \capa{9}{\node[punto,label={[etiqueta]right:$C'$}] at (4.6748,3.9834) {};}
  % 10: el núcleo se encoge hacia E al replicar la economía
  \capa{10}{\draw[flecha] plot[domain=3.2:3.78,samples=20] (\x,{cc(\x)});}
  \capa{10}{\draw[flecha] plot[domain=4.6:4.04,samples=20] (\x,{cc(\x)});}
  \capa{10}{\node[etiqueta,text=tceerojo] at (3.5,2.45) {$r\to\infty$};}
\end{tikzpicture}
\end{document}
 muestra solo
% esas capas (en el vídeo, \tceepaso sigue mandando).
\makeatletter
\ifdefined\cajacapas
  \let\capamaestro\capa
  \renewcommand{\capa}[2]{\edef\cajaesta{\noexpand\in@{,#1,}{,\cajacapas,}}\cajaesta
    \ifin@\capamaestro{#1}{#2}\fi}
\fi
\makeatother
\begin{document}
\begin{tikzpicture}[tcee,x=1cm,y=1.1cm,
    declare function={
      ua(\x,\u)=(\u/\x^0.4)^(1/0.6);
      ub(\x,\v)=6-(\v/(10-\x)^0.6)^(1/0.4);
      cc(\x)=9*\x/((2/3)*(10-\x)+1.5*\x);
      odax(\p)=0.4*(1.5*\p+5)/\p;  oday(\p)=0.6*(1.5*\p+5);
      odbx(\p)=10-0.6*(8.5*\p+1)/\p; odby(\p)=6-0.4*(8.5*\p+1);},
    llave/.style={decorate,decoration={brace,amplitude=5pt},thin,tceegris}]
  % 1: la caja, sus orígenes, los ejes de cada consumidor y la dotación
  \capa{1}{
    \draw[thick] (0,0) rectangle (10,6);
    \node[below left,text=tceemorado] at (0,0) {$O_A$};
    \node[above right,text=tceeverde] at (10,6) {$O_B$};
    \draw[->,thin,tceemorado] (0.4,-0.3) -- (1.4,-0.3) node[right,etiqueta] {$X_A$};
    \draw[->,thin,tceemorado] (-0.3,0.4) -- (-0.3,1.4) node[above,etiqueta] {$Y_A$};
    \draw[->,thin,tceeverde] (9.6,6.3) -- (8.6,6.3) node[left,etiqueta] {$X_B$};
    \draw[->,thin,tceeverde] (10.3,5.6) -- (10.3,4.6) node[below,etiqueta] {$Y_B$};
    \draw[llave] (10,-0.75) -- (0,-0.75) node[midway,below=5pt,text=black] {$\bar X$};
    \draw[llave] (-0.95,0) -- (-0.95,6) node[midway,left=5pt,text=black] {$\bar Y$};
  }
  \capa{1}{\node[punto,label={[etiqueta]below left:$\vect{\omega}$}] (w) at (1.5,5) {};}
  % 2: curvas de indiferencia que pasan por ω y la lente
  \capa{2}{\draw[curva1] plot[domain=1.2:9.3,samples=120] (\x,{ua(\x,3.09017)})
    node[etiqueta,right] {$U_A^0$};}
  \capa{2}{\draw[curva3] plot[domain=0.2:7.38,samples=120] (\x,{ub(\x,3.60986)})
    node[etiqueta,right] {$U_B^0$};}
  \capa{2}{\fill[area,fill opacity=0.12] plot[domain=1.5:6.7133,samples=100] (\x,{ua(\x,3.09017)})
    -- plot[domain=6.7133:1.5,samples=100] (\x,{ub(\x,3.60986)}) -- cycle;}
  % 3: precios de desequilibrio (p_X/p_Y = 0,23): lo que elige cada uno y los excesos
  \capa{3}{\draw[curva4,desplazada] (0.5,5.23) -- (9.6,3.137) node[etiqueta,anchor=north east] {$p_X/p_Y$};}
  \capa{3}{\node[punto,label={[etiqueta]above left:$\vect{x}_A$}] at (9.2957,3.207) {};}
  \capa{3}{\node[punto,label={[etiqueta,label distance=3pt]above:$\vect{x}_B$}] at (2.2913,4.818) {};}
  \capa{3}{\draw[curvanueva] (2.2913,4.818) -- (9.2957,4.818) node[midway,above,etiqueta] {$z_X>0$};}
  \capa{3}{\draw[curvanueva] (9.2957,4.818) -- (9.2957,3.207) node[pos=0.3,left,etiqueta] {$z_Y<0$};}
  % 4: curva de oferta-demanda de A
  \capa{4}{\draw[curva1,desplazada,thick] plot[domain=2.2222:0.23,samples=150,variable=\p]
    ({odax(\p)},{oday(\p)});}
  \capa{4}{\node[etiqueta,text=tceemorado,above] at (6.66,3.297) {$\mathrm{OD}_A$};}
  % 5: curva de oferta-demanda de B
  \capa{5}{\draw[curva3,desplazada,thick] plot[domain=0.17647:1.25,samples=150,variable=\p]
    ({odbx(\p)},{odby(\p)}) node[etiqueta,below] {$\mathrm{OD}_B$};}
  % 6: equilibrio E, recta de precios de equilibrio y curvas tangentes en E
  \capa{6}{\node[punto,label={[etiqueta]above:$E$}] (E) at (3.9077,3.5442) {};}
  \capa{6}{\draw[curva4] (0.6,5.5442) -- (9.7686,0) node[etiqueta,below] {$p^*_X/p^*_Y$};}
  \capa{6}{\draw[curva1] plot[domain=1.9:9.3,samples=120] (\x,{ua(\x,3.68534)})
    node[etiqueta,right] {$U_A^*$};}
  \capa{6}{\draw[curva3] plot[domain=0.2:6.6,samples=120] (\x,{ub(\x,4.23592)})
    node[etiqueta,right] {$U_B^*$};}
  % 7: curva de contrato (1TFEB: E está en ella)
  \capa{7}{\draw[curva2] plot[domain=0:10,samples=150] (\x,{cc(\x)});}
  \capa{7}{\node[etiqueta,text=tceeazul,anchor=east] at (7.9,5.75) {curva de contrato};}
  % 8: 2TFEB: otro óptimo E', su recta de precios y la redistribución ω → ω'
  \capa{8}{\node[punto,label={[etiqueta]right:$E'$}] at (2.2,2.3294) {};}
  \capa{8}{\draw[curva4] (0.9,3.2471) -- (5.5,0) ;}
  \capa{8}{\node[punto,label={[etiqueta]left:$\vect{\omega}'$}] (w2) at (1.5,2.8235) {};}
  \capa{8}{\draw[flecha,shorten <=3pt,shorten >=3pt] (w) -- (w2);}
  % 9: Edgeworth: el tramo CC' de la curva de contrato dentro de la lente
  \capa{9}{\draw[curva2,line width=5pt,opacity=0.35] plot[domain=3.1484:4.6748,samples=40] (\x,{cc(\x)});}
  \capa{9}{\node[punto,label={[etiqueta]left:$C$}] at (3.1484,3.0500) {};}
  \capa{9}{\node[punto,label={[etiqueta]right:$C'$}] at (4.6748,3.9834) {};}
  % 10: el núcleo se encoge hacia E al replicar la economía
  \capa{10}{\draw[flecha] plot[domain=3.2:3.78,samples=20] (\x,{cc(\x)});}
  \capa{10}{\draw[flecha] plot[domain=4.6:4.04,samples=20] (\x,{cc(\x)});}
  \capa{10}{\node[etiqueta,text=tceerojo] at (3.5,2.45) {$r\to\infty$};}
\end{tikzpicture}
\end{document}
